\chapter{Custom Gradients for LEDH-PFPF-OT}
\label{ch:bf-ledh-pfpf-ot-custom-gradient}

The preceding chapters introduced three ingredients separately: local exact
Daum--Huang particle flow, proposal-corrected PF-PF weighting, and
entropy-regularized OT resampling.  In the OT sequence immediately before this
chapter, the deterministic baseline fixed the transport-object and barycentric-output
contract, while the entropic chapter fixed the exact teacher object, the finite
Sinkhorn route that approximates it, and the rule that later differentiation claims
must attach to that executed finite route.  This chapter explains how those pieces can
be differentiated without asking an automatic-differentiation system to retain
the entire filtering graph in memory.  The key object is a custom
vector-Jacobian product, or VJP, for the executed LEDH-PFPF-OT scalar.

The chapter is intentionally narrow.  It does not claim that the relaxed
OT-resampled particle filter is the same object as categorical resampling.  It
does not claim that finite Sinkhorn iteration is the exact unregularized OT
optimizer.  It does not claim HMC correctness for a nonlinear model.  Its claim
is more basic and more useful: if the value path is a declared finite
LEDH-PFPF-OT computation, then the gradient path must be the derivative of that
same computation, and the expensive OT block has a structured adjoint.

\section{The scalar whose derivative is wanted}
\label{sec:bf-ledh-ot-scalar}

Let $\theta$ denote the model parameter and let all random numbers used by the
diagnostic or sampler target be fixed as part of the deterministic value path.
For one filtering pass, write the executed scalar as
\begin{equation}
\label{eq:bf-ledh-ot-executed-scalar}
  \mathcal L_K(\theta)
  =
  \operatorname{Reduce}
  \left(
    \operatorname{LEDH\mbox{-}PFPF\mbox{-}OT}_K
    (\theta; y_{1:T}, \xi)
  \right),
\end{equation}
where $\xi$ denotes the fixed particle and process-noise streams, $K$ denotes
the finite Sinkhorn or annealed-transport iteration budget, and
\(\operatorname{Reduce}\) is the declared scalar reduction, such as a mean over
fixed seeds.  The derivative needed by HMC or optimization is
\[
  \nabla_\theta \mathcal L_K(\theta),
\]
not the derivative of a different resampling rule, a different Sinkhorn stopping
policy, or a different random seed construction.

\begin{definition}[Same-scalar custom-gradient contract]
\label{def:bf-ledh-ot-same-scalar}
A custom-gradient implementation for LEDH-PFPF-OT is same-scalar if its primal
return is exactly \(\mathcal L_K(\theta)\) and its VJP returns
\[
  v^\top D_\theta \mathcal L_K(\theta)
\]
for every upstream scalar multiplier \(v\).  For a scalar log target the usual
case is \(v=1\).
\end{definition}

This contract is the same HMC principle as Chapter~\ref{ch:bf-custom-gradient-wrappers},
but the present chapter specializes it to the LEDH-PFPF-OT computation.  The
particle-flow part follows the Li--Coates invertible-particle-flow proposal
logic \citep{li2017particle}; the OT relaxation follows the entropy-regularized
resampling route of \citet{corenflos2021differentiable}, with the Sinkhorn
background of \citet{cuturi2013sinkhorn} and
\citet{peyre2019computational}.
The derivative identities below are project derivations for the declared
BayesFilter finite program; the citations explain the particle-flow,
entropy-regularized OT, and Sinkhorn objects being differentiated.

\section{Forward decomposition}
\label{sec:bf-ledh-ot-forward-decomposition}

At one observation time, the canonical value path is transition-first and can
be written as a composition
\begin{equation}
\label{eq:bf-ledh-ot-forward-composition}
  \theta
  \longmapsto
  (X^-, \log w^-, A_{\mathrm{flow}})
  \longmapsto
  (X, \log w)
  \longmapsto
  (\Delta_t,X^\star,\log w^\star)
  \longmapsto
  \mathcal L_t .
\end{equation}
Here \(X^-\in\mathbb R^{N\times d}\) is the pre-observation proposal cloud,
\(\log w^-\) are proposal-correction log weights, and
\(A_{\mathrm{flow}}\) denotes the LEDH auxiliary quantities needed for the
proposal density and log-determinant correction.  The local LEDH construction
uses particle-local linearizations and particle-local covariance state as in
Li--Coates PF-PF.  Corrected logits first produce the current likelihood
increment \(\Delta_t\) and normalized weights.  Only after that increment has
entered the scalar does an active reset map the weighted cloud to an
equal-weight cloud
\[
  (X,\log w) \longmapsto X^\star.
\]

This chapter focuses on the last map because it is the part that can dominate
memory.  Naive reverse-mode autodiff through all pairwise costs, Sinkhorn
updates, and barycentric products may retain objects of size \(O(KN^2)\).  A
custom VJP instead treats the reset block as a finite differentiable program and
applies its adjoint by recomputation, streaming reductions, or checkpointing.

The canonical reset in this chapter is Contract E--Chol.  A raw barycentric
cloud, compact forward sensitivity of that raw scalar, or full-history/manual
reverse derivative of that raw scalar is historical diagnostic evidence only.
Even an algebraically correct derivative of the raw-barycentric scalar is wrong
relative to the canonical Contract E total-gradient claim.  There is no
canonical-to-raw fallback.

\section{Canonical row quotient and Contract E composition}
\label{sec:bf-ledh-contract-e-total-composition}

The executed finite streaming transport does not get to assume its intended
row marginal is exact.  It emits
\begin{equation}
\label{eq:bf-ledh-contract-e-row-quotient}
  M_i=\sum_jP_{ij},\qquad
  Q_i=\sum_jP_{ij}X_j,\qquad
  Y_i^+=Q_i/M_i .
\end{equation}
The canonical route has no denominator floor.  For a tangent and an upstream
cotangent,
\begin{equation}
\label{eq:bf-ledh-contract-e-row-quotient-jvp-vjp}
  \dot Y_i^+=\frac{\dot Q_i-Y_i^+\dot M_i}{M_i},\qquad
  G_i^Q=G_i^{Y^+}/M_i,\qquad
  G_i^M=-\langle G_i^{Y^+},Y_i^+\rangle/M_i .
\end{equation}
Both \(G^Q\) and \(G^M\) must enter the finite Sinkhorn pullback.

Contract E consumes the same source cloud \(X\), normalized probabilities
\(w\), row-quotient cloud \(Y^+\), fixed residual design \(\Xi\), and prepared
ridge \(\lambda\).  It depends directly on the weighted source moments as well
as indirectly through transport.  In row-vector storage its local moments and
reset are
\begin{align}
  \mu_w&=\sum_iw_iX_i,&
  \Sigma_w&=\sum_iw_i(X_i-\mu_w)(X_i-\mu_w)^\top,\\
  \bar Y^+&=N^{-1}\sum_iY_i^+,&
  \Sigma_+&=N^{-1}\sum_i(Y_i^+-\bar Y^+)(Y_i^+-\bar Y^+)^\top,
\end{align}
\begin{equation}
\label{eq:bf-ledh-contract-e-affine-reset}
  G_+=\operatorname{sym}(\Sigma_w-\Sigma_+),\quad
  B_\lambda=\operatorname{chol}(G_++\lambda I),\quad
  \widetilde Y=Y^++\Xi B_\lambda^\top,
\end{equation}
\begin{equation}
  \widetilde\Sigma=\frac1N\sum_i
    (\widetilde Y_i-\bar{\widetilde Y})
    (\widetilde Y_i-\bar{\widetilde Y})^\top,\quad
  A_\lambda=
    \operatorname{chol}(\Sigma_w+\lambda I)
    \operatorname{chol}(\widetilde\Sigma+\lambda I)^{-1},
\end{equation}
\begin{equation}
  X_i^\star=\mu_w+
    (\widetilde Y_i-\bar{\widetilde Y})A_\lambda^\top.
\end{equation}
The implementation evaluates the last operator by a triangular solve rather
than forming an explicit inverse.  Its cloud VJP first returns
\[
  G_X^{\mathrm{mom}},\quad G_w^{\mathrm{mom}},\quad G_{Y^+},
  \quad G_\Xi,\quad g^\lambda .
\]
Composing \(G_{Y^+}\) through
\eqref{eq:bf-ledh-contract-e-row-quotient-jvp-vjp} and the finite Sinkhorn
program returns transport terms.  In common coordinates the canonical totals
are
\begin{equation}
\label{eq:bf-ledh-contract-e-total-source-pullback}
  G_X=G_X^{\mathrm{mom}}+G_X^{\mathrm{transport}},\qquad
  G_w=G_w^{\mathrm{mom}}+G_w^{\mathrm{transport}} .
\end{equation}
The implementation naturally receives a normalized-log-weight cotangent from
transport.  It therefore forms
\begin{equation}
\label{eq:bf-ledh-contract-e-total-logweight-pullback}
  G_{\log w}
  =G_{\log w}^{\mathrm{transport}}+w\odot G_w^{\mathrm{mom}}
\end{equation}
before applying the log-normalization pullback exactly once.  A
probability-coordinate moment adjoint must not be added directly to a log-weight
adjoint, and normalization must not be applied twice.  The residual and ridge
adjoints are returned for parity, but the canonical route fixes their prepared
values, so \(\dot\Xi=0\) and \(\dot\lambda=0\).  A future
parameter-dependent design or ridge would be a different route and would have
to differentiate that dependence fully.

\section{Transport orientation}
\label{sec:bf-ledh-ot-orientation}

The OT chapters used a coupling \(\Pi\) with source index first and output slot
second.  The TensorFlow LEDH-PFPF-OT route is more naturally described by the
row-stochastic transport matrix
\begin{equation}
\label{eq:bf-ledh-ot-row-transport}
  T = N\Pi^\top,
  \qquad
  T\mathbf 1=\mathbf 1,
  \qquad
  T^\top\mathbf 1=Nw,
\end{equation}
where \(w=\operatorname{softmax}(\log w)\).  The transported equal-weight cloud
is
\begin{equation}
\label{eq:bf-ledh-ot-barycentric-output}
  \widetilde X = T X,
  \qquad
  \widetilde x_j=\sum_{i=1}^N T_{ji}x_i .
\end{equation}
Thus row \(j\) of \(T\) tells how output slot \(j\) averages the source
particles.  The row sums make each output particle a barycenter.  The column
sums say that source particle \(i\) contributes total mass proportional to its
weight.

\begin{remark}[Why orientation matters]
\label{rem:bf-ledh-ot-orientation}
If a derivation uses \(\Pi\) but the code applies \(T X\), transpose errors are
easy to introduce.  The VJP formulas below use \(T\), not \(\Pi\).  A reader who
wants the coupling form can substitute \(T=N\Pi^\top\).
\end{remark}

\section{VJP of the barycentric product}
\label{sec:bf-ledh-ot-barycentric-vjp}

The simplest and most important adjoint identity is the VJP of the barycentric
matrix product.

\begin{proposition}[Barycentric VJP]
\label{prop:bf-ledh-ot-barycentric-vjp}
Let \(\widetilde X=T X\), with \(T\in\mathbb R^{N\times N}\) and
\(X\in\mathbb R^{N\times d}\).  Use the Frobenius inner product for matrix
adjoints, and let \(G_{\widetilde X}\in\mathbb R^{N\times d}\) be the upstream
adjoint for \(\widetilde X\).  Then the direct adjoints are
\begin{equation}
\label{eq:bf-ledh-ot-vjp-barycentric}
  G_X^{\mathrm{direct}} = T^\top G_{\widetilde X},
  \qquad
  G_T = G_{\widetilde X}X^\top .
\end{equation}
\end{proposition}

\begin{proof}
Taking a differential gives
\[
  d\widetilde X = dT\,X + T\,dX .
\]
Pairing with \(G_{\widetilde X}\) under the Frobenius inner product,
\begin{align*}
  \langle G_{\widetilde X},d\widetilde X\rangle
  &=
  \langle G_{\widetilde X},dT\,X\rangle
  +
  \langle G_{\widetilde X},T\,dX\rangle  \\
  &=
  \langle G_{\widetilde X}X^\top,dT\rangle
  +
  \langle T^\top G_{\widetilde X},dX\rangle .
\end{align*}
The coefficients of \(dT\) and \(dX\) are exactly
\eqref{eq:bf-ledh-ot-vjp-barycentric}.
\end{proof}

Equation~\eqref{eq:bf-ledh-ot-vjp-barycentric} is only the direct part of the
particle adjoint.  Because \(T\) itself depends on \(X\) and \(\log w\), the
complete VJP also pulls \(G_T\) backward through the Sinkhorn or
annealed-transport construction.

\section{A normalized transport matrix}
\label{sec:bf-ledh-ot-normalized-transport}

For fixed final potentials \(f,g\in\mathbb R^N\), fixed normalized log weights
\(\log w_i\), and cost \(C_{ji}=c(z_j,z_i)\), the executed row-stochastic
transport can be written in column-normalized form
\begin{equation}
\label{eq:bf-ledh-ot-transport-form}
  T_{ji}
  =
  N w_i
  \frac{
    \exp\{(f_j+g_i-C_{ji})/\varepsilon\}
  }{
    \sum_{r=1}^N
    \exp\{(f_r+g_i-C_{ri})/\varepsilon\}
  } .
\end{equation}
This expression enforces \(T^\top\mathbf 1=Nw\) exactly for the displayed
normalization.  The row condition \(T\mathbf 1\approx\mathbf 1\) is enforced by
the potentials produced by finite Sinkhorn or annealed Sinkhorn iteration.
In the TensorFlow implementation this is the convention used by the final
transport application: \(f\) indexes output rows, \(g\) indexes source columns,
the normalization is a column log-sum-exp over output rows, and the transported
particles are obtained by the matrix product \(T X\).

Define
\begin{equation}
\label{eq:bf-ledh-ot-column-prob}
  p_{ji}
  =
  \frac{T_{ji}}{Nw_i},
  \qquad
  \sum_{j=1}^N p_{ji}=1.
\end{equation}
The vector \(p_{\cdot i}\) is the normalized column distribution associated with
source particle \(i\).

\begin{proposition}[VJP of the normalized transport output]
\label{prop:bf-ledh-ot-normalized-transport-vjp}
At fixed finite values of \(f,g\in\mathbb R^N\),
\(C\in\mathbb R^{N\times N}\), \(\varepsilon>0\), and normalized weights
\(w_i>0\), \(\sum_i w_i=1\), let \(T\) be defined by
\eqref{eq:bf-ledh-ot-transport-form}.  Let
\(G_T\in\mathbb R^{N\times N}\) be the upstream adjoint for \(T\).  For each
source column \(i\), set
\begin{equation}
\label{eq:bf-ledh-ot-column-mean-adjoint}
  \bar G_i=\sum_{r=1}^N p_{ri}(G_T)_{ri},
  \qquad
  S_{ji}=T_{ji}\bigl((G_T)_{ji}-\bar G_i\bigr).
\end{equation}
Then the VJP of \eqref{eq:bf-ledh-ot-transport-form} satisfies
\begin{align}
\label{eq:bf-ledh-ot-vjp-logw}
  (G_{\log w})_i
  &\mathrel{+}=
  \sum_{j=1}^N (G_T)_{ji}T_{ji},\\
\label{eq:bf-ledh-ot-vjp-f}
  (G_f)_j
  &\mathrel{+}=
  \frac{1}{\varepsilon}\sum_{i=1}^N S_{ji},\\
\label{eq:bf-ledh-ot-vjp-g}
  (G_g)_i
  &\mathrel{+}=
  \frac{1}{\varepsilon}\sum_{j=1}^N S_{ji},\\
\label{eq:bf-ledh-ot-vjp-cost}
  (G_C)_{ji}
  &\mathrel{+}=
  -\frac{1}{\varepsilon}S_{ji}.
\end{align}
If \(\log w\) is itself produced by a \(\operatorname{logsoftmax}\), then
\eqref{eq:bf-ledh-ot-vjp-logw} must be composed with the
\(\operatorname{logsoftmax}\) VJP.
\end{proposition}

\begin{proof}
Let
\[
  a_{ji}=\frac{f_j+g_i-C_{ji}}{\varepsilon},
  \qquad
  p_{ji}=\frac{\exp(a_{ji})}{\sum_r\exp(a_{ri})},
  \qquad
  T_{ji}=Nw_i p_{ji}.
\]
For a perturbation of one column \(i\),
\[
  dT_{ji}
  =
  T_{ji}
  \left[
    d\log w_i
    + da_{ji}
    - \sum_{r=1}^N p_{ri}da_{ri}
  \right].
\]
Pairing this expression with \((G_T)_{ji}\) and collecting the coefficient of
\(da_{ji}\) gives \(S_{ji}=T_{ji}((G_T)_{ji}-\bar G_i)\).  Since
\[
  da_{ji}
  =
  \frac{1}{\varepsilon}
  (df_j+dg_i-dC_{ji}),
\]
the adjoints for \(f\), \(g\), and \(C\) are
\eqref{eq:bf-ledh-ot-vjp-f}--\eqref{eq:bf-ledh-ot-vjp-cost}.  The coefficient
of \(d\log w_i\) is \(\sum_j (G_T)_{ji}T_{ji}\), giving
\eqref{eq:bf-ledh-ot-vjp-logw}.
\end{proof}

This proposition is the core OT VJP product.  It says that the upstream matrix
\(G_T\) should first be column-centered under the transport-induced probability
\(p_{\cdot i}\).  Only after that centering should the adjoint be pushed into
the potentials and costs.  In the displayed column-normalized map,
\(\sum_j S_{ji}=0\), so the final-potential adjoint \(G_g\) is identically zero
at this last normalization stage.  This is not a bug: adding the same \(g_i\)
to every logit in one normalized column does not change that column's
probabilities.  Earlier Sinkhorn-loop adjoints may still involve the variables
that produced \(g\).

\section{VJP of the quadratic cost}
\label{sec:bf-ledh-ot-cost-vjp}

For the common squared-distance cost
\[
  C_{ji}=\frac{1}{2}\|z_j-z_i\|^2,
\]
where the same scaled particle cloud \(Z\) supplies the row and column points,
the cost adjoint has another simple form.

\begin{proposition}[Quadratic-cost VJP]
\label{prop:bf-ledh-ot-cost-vjp}
Let \(G_C\) be the upstream adjoint for
\(C_{ji}=\frac12\|z_j-z_i\|^2\), where
\(Z=(z_1,\ldots,z_N)^\top\in\mathbb R^{N\times d}\).  The contribution to the
adjoint of \(z_k\in\mathbb R^d\) is
\begin{equation}
\label{eq:bf-ledh-ot-vjp-cost-points}
  (G_Z)_k
  \mathrel{+}=
  \sum_{i=1}^N (G_C)_{ki}(z_k-z_i)
  -
  \sum_{j=1}^N (G_C)_{jk}(z_j-z_k).
\end{equation}
\end{proposition}

\begin{proof}
The differential of one cost entry is
\[
  dC_{ji}=(z_j-z_i)^\top(dz_j-dz_i).
\]
Multiplying by \((G_C)_{ji}\), summing over \(j,i\), and collecting the
coefficient of \(dz_k\) gives the two sums in
\eqref{eq:bf-ledh-ot-vjp-cost-points}: the first from terms where \(k\) is the
row index, and the second from terms where \(k\) is the column index.
\end{proof}

If centering, scaling, or key construction is stopped by design, the VJP stops
there too.  If the selected mode differentiates through centering or scaling,
then \eqref{eq:bf-ledh-ot-vjp-cost-points} must be further composed with the
VJP of those preprocessing maps.

\section{Softmin and Sinkhorn adjoints}
\label{sec:bf-ledh-ot-softmin-vjp}

The finite Sinkhorn or annealed-transport step is a loop of softmin operations.
The exact schedule can differ between a fixed-target Sinkhorn route and the
FilterFlow-style annealed route, but the local adjoint primitive is the same.

\begin{definition}[Softmin block]
\label{def:bf-ledh-ot-softmin}
For a query index \(j\), costs \(C_{jr}\), values \(v_r\), and temperature
\(\varepsilon>0\), define
\begin{equation}
\label{eq:bf-ledh-ot-softmin}
  m_j
  =
  -\varepsilon
  \log
  \sum_{r=1}^N
  \exp\{(v_r-C_{jr})/\varepsilon\}.
\end{equation}
\end{definition}

\begin{proposition}[Softmin VJP for fixed temperature]
\label{prop:bf-ledh-ot-softmin-vjp}
Assume \(\varepsilon>0\) is held fixed and all displayed exponentials are
finite.  Let \(G_m\in\mathbb R^N\) be the upstream adjoint for \(m\), and define
\[
  q_{jr}
  =
  \frac{\exp\{(v_r-C_{jr})/\varepsilon\}}
       {\sum_s \exp\{(v_s-C_{js})/\varepsilon\}}.
\]
For fixed \(\varepsilon\),
\begin{align}
\label{eq:bf-ledh-ot-softmin-vjp-values}
  (G_v)_r
  &\mathrel{+}=
  -\sum_{j=1}^N (G_m)_j q_{jr},\\
\label{eq:bf-ledh-ot-softmin-vjp-costs}
  (G_C)_{jr}
  &\mathrel{+}=
  (G_m)_j q_{jr}.
\end{align}
\end{proposition}

\begin{proof}
Differentiating \eqref{eq:bf-ledh-ot-softmin} gives
\[
  dm_j
  =
  -\sum_r q_{jr}(dv_r-dC_{jr})
  =
  \sum_r q_{jr}dC_{jr}-\sum_r q_{jr}dv_r.
\]
Pairing with \(G_m\) and collecting coefficients yields
\eqref{eq:bf-ledh-ot-softmin-vjp-values} and
\eqref{eq:bf-ledh-ot-softmin-vjp-costs}.
\end{proof}

The sign and scale in Proposition~\ref{prop:bf-ledh-ot-softmin-vjp} are for a
potential-scale value \(v\).  Some implementations, including the TensorFlow
helper used in the current BayesFilter route, call the same operation with a
log-scale input \(h_r\):
\[
  m_j=-\varepsilon\log\sum_r
  \exp\{h_r-C_{jr}/\varepsilon\}.
\]
This is the same formula with \(v_r=\varepsilon h_r\).  Therefore the local VJP
with respect to \(h_r\) is \((G_h)_r\mathrel{+}=
-\varepsilon\sum_j(G_m)_j q_{jr}\), while the VJP with respect to a
potential-scale component \(v_r\) remains
\eqref{eq:bf-ledh-ot-softmin-vjp-values}.  If \(h_r=\log\alpha_r+b_r/\varepsilon\),
then the adjoint to \(b_r\) again has the potential-scale coefficient
\(-\sum_j(G_m)_j q_{jr}\).

For an annealed schedule, the loop state contains several potentials, for
example \((a_y,b_x,a_x,b_y)\).  Abstractly write the finite loop as
\begin{equation}
\label{eq:bf-ledh-ot-loop-state}
  s_{k+1}=F_k(s_k;Z,\log w,\varepsilon_k),
  \qquad k=0,\ldots,K-1.
\end{equation}
The reverse loop is then
\begin{equation}
\label{eq:bf-ledh-ot-reverse-loop}
  (G_{s_k},G_Z,G_{\log w})
  \mathrel{+}=
  D F_k(s_k;Z,\log w,\varepsilon_k)^\top G_{s_{k+1}},
  \qquad k=K-1,\ldots,0.
\end{equation}
Each multiplication by \(D F_k^\top\) is implemented by applying the softmin
VJP above, the averaging VJP for damped updates, and the cost VJP for the
pairwise squared distances.  This is ordinary reverse-mode differentiation of a
finite program, but expressed as an explicit adjoint scan rather than as a raw
tape over the whole program.

\begin{proposition}[Finite-Sinkhorn VJP by reverse scan]
\label{prop:bf-ledh-ot-finite-sinkhorn-vjp}
Assume the finite transport block is the deterministic program
\[
  (T,s_K)=\mathcal S_K(Z,\log w)
\]
obtained by a fixed number of differentiable softmin, averaging, cost, and
normalization operations with fixed branch decisions.  Then a custom VJP for
\(\mathcal S_K\) is obtained by:
\begin{enumerate}[label=(\roman*)]
  \item applying Proposition~\ref{prop:bf-ledh-ot-normalized-transport-vjp} to
  move \(G_T\) into final potentials, costs, and log weights;
  \item applying Proposition~\ref{prop:bf-ledh-ot-cost-vjp} to move final-cost
  adjoints into \(Z\);
  \item scanning \eqref{eq:bf-ledh-ot-reverse-loop} backward through the finite
  potential updates;
  \item composing any resulting \(G_Z\) with the chosen centering, scaling, and
  stop-gradient policy.
\end{enumerate}
The result equals the VJP of the executed finite program.
\end{proposition}

\begin{proof}
The finite transport block is a composition of differentiable maps under the
fixed branch and fixed iteration assumptions.  Propositions
\ref{prop:bf-ledh-ot-normalized-transport-vjp},
\ref{prop:bf-ledh-ot-cost-vjp}, and
\ref{prop:bf-ledh-ot-softmin-vjp} give the VJP for the nontrivial primitive
maps.  The VJP of a composition is obtained by applying the transposed local
Jacobians in reverse order.  That reverse composition is exactly the scan
described above.
\end{proof}

\section{Why a custom VJP is memory efficient}
\label{sec:bf-ledh-ot-memory}

The naive reverse graph for the OT block stores pairwise costs, logits,
normalizers, potentials, and intermediate loop states for every time step.  With
\(N\) particles and \(K\) Sinkhorn updates, this can scale like \(O(KN^2)\) in
stored tensors before the rest of the filtering graph is considered.

The custom VJP changes the storage problem.  The primal path must retain or
reconstruct only enough information to replay the finite program:
\begin{itemize}
  \item the input particles and log weights entering the OT block;
  \item the branch choices, masks, iteration count, temperature schedule, and
  stabilization policy;
  \item final potentials, and optionally sparse checkpoints of intermediate
  potentials;
  \item any stopped quantities needed to reproduce the executed scalar.
\end{itemize}
During the backward pass, pairwise blocks can be recomputed and reduced
immediately.  This trades additional arithmetic for lower peak memory.  The
trade is natural in HMC, because one scalar log target needs one reverse VJP to
obtain all parameter-score components, while a forward-mode JVP would need one
pass per requested direction.

\section{Same-cloud geometry reuse in the finite JVP}
\label{sec:bf-ledh-ot-same-cloud-cache}

The canonical finite forward-JVP route repeatedly evaluates softmin operations
on the same scaled particle cloud during one OT solve.  Let
\(Z=(z_1,\ldots,z_N)\), let the parameter tangent axis have arbitrary size
\(P\), and define
\begin{align}
\label{eq:bf-ledh-cached-cost}
  C_{ij}
  &=\frac12\lVert z_i-z_j\rVert^2,\\
\label{eq:bf-ledh-cached-cost-jvp}
  \dot C_{ij}^{(p)}
  &=(z_i-z_j)^\top
    \bigl(\dot z_i^{(p)}-\dot z_j^{(p)}\bigr),
  \qquad p=1,\ldots,P.
\end{align}
The Sinkhorn potentials change from iteration to iteration, but \(Z\) and its
tangent do not change within that finite solve.  Therefore the same \(C\) and
\(\dot C\) appear in every annealed softmin, every terminal-balancing softmin,
and the final transport assembly.

\begin{proposition}[Same-cloud cache preserves the finite JVP]
\label{prop:bf-ledh-same-cloud-cache}
Fix the input cloud \(Z\), its parameter tangents \(\dot Z\), the finite
temperature schedule, the potential initialization, the iteration counts, and
all branch decisions for one OT solve.  Replacing each recomputation of
\(C(Z)\) and \(D C(Z)[\dot Z]\) by the cached tensors in
\eqref{eq:bf-ledh-cached-cost}--\eqref{eq:bf-ledh-cached-cost-jvp} leaves the
mathematical softmin values and JVPs unchanged.  Consequently it leaves the
declared finite Sinkhorn, terminal-balance, row-quotient, and Contract E total
JVP unchanged in exact arithmetic.
\end{proposition}

\begin{proof}
For log-scale inputs \(h_j\), one softmin and its JVP are
\begin{align*}
  m_i
  &=-\varepsilon\log\sum_j
    \exp\{h_j-C_{ij}/\varepsilon\},\\
  \dot m_i
  &=-\dot\varepsilon\log\sum_j e^{\ell_{ij}}
    -\varepsilon\sum_j q_{ij}\dot\ell_{ij},\\
  \dot\ell_{ij}
  &=\dot h_j-\frac{\dot C_{ij}}{\varepsilon}
    +\frac{C_{ij}\dot\varepsilon}{\varepsilon^2},
  \qquad
  q_{ij}=\frac{e^{\ell_{ij}}}{\sum_r e^{\ell_{ir}}}.
\end{align*}
Within the fixed solve, every occurrence of \(C_{ij}\) is the deterministic
function in \eqref{eq:bf-ledh-cached-cost}, and every occurrence of
\(\dot C_{ij}\) is its directional derivative in
\eqref{eq:bf-ledh-cached-cost-jvp}.  Substituting stored copies of those same
tensors changes neither displayed expression.  The finite solver is a
composition of these softmins, affine potential updates, reductions, and the
final transport and Contract E maps.  Induction over the fixed iteration
sequence therefore gives the same value and tangent at every subsequent state.
\end{proof}

This optimization addresses a concrete generic bottleneck.  With \(K\)
annealed Sinkhorn updates and \(L\) terminal-balance updates, the current
same-cloud JVP evaluates four softmins before the loop, four per Sinkhorn
update, two after it, and two per balance update.  Thus it otherwise constructs
or differentiates pairwise geometry
\[
  4K+2L+6
\]
times per OT solve, before counting the final transport assembly.  For the
measured \(K=20,L=8\) witness this is \(102\) repeated large geometry/JVP
traversals per filtering time step.  Caching reduces those geometry
constructions to one \(C\) and one \(\dot C\) per solve; it does not reduce the
required softmax reductions or change either iteration count.

The storage tradeoff is explicit.  For batch size \(B\), particle count \(N\),
parameter count \(P\), and scalar width \(s\) bytes, the two cached tensors
require approximately
\begin{equation}
\label{eq:bf-ledh-cached-memory}
  s B N^2(1+P)
\end{equation}
bytes before compiler liveness and other filter state are counted.  The formula
supports any positive \(P\), but does not imply that dense caching is feasible
for every \((B,N,P)\).  The cache is therefore a scoped optimization: it may be
enabled only when the route uses one exact same-cloud tile and a scope-specific
GPU memory/performance check passes.  Otherwise the streamed recomputation path
remains the correct fallback.  No cache may cross a time step, parameter
evaluation, changed cloud, changed tangent, or changed dtype.

On the July 2026 TF32/XLA diagnostic witness with \(T=50\), \(N=1024\),
\(K=20\), and \(L=8\), five synchronized warm calls at \(B=1\) had median
times \(3.013\) seconds without the cache and \(0.302\) seconds with it, a
descriptive \(9.97\)-fold speedup.  At \(B=16\), three warm calls had medians
\(25.922\) and \(2.591\) seconds, respectively, a descriptive
\(10.00\)-fold speedup.  Both arms used one XLA \texttt{StatelessWhile}, the
repository \(K=N\) chunk policy, the same prepared inputs, and the same 8 GiB
device cap.  The relative score-norm differences were \(3.24\times10^{-5}\)
and \(1.78\times10^{-5}\), respectively; all finite, replay, chart, marginal,
and reset gates passed.  These differences are floating-point evaluation-order
effects, not a change in the exact-arithmetic finite target.  LGSSM is only the
measured performance witness: the implementation and parity tests carry an
arbitrary final parameter axis and do not use Kalman or LGSSM formulas.

An exact parameter-axis batching candidate for the Contract E Cholesky reset
was also tested.  Although it matched the per-direction algebra on
\(P\in\{1,3,5,18\}\), its measured \(B=1\) median rose from \(0.302\) to
\(0.413\) seconds.  It is therefore rejected as an optimization and the
per-direction reset JVP remains the active route.  Correct vectorization is not
evidence of a performance gain.

\section{Embedding the OT VJP in the filtering loop}
\label{sec:bf-ledh-ot-filter-loop}

The full LEDH-PFPF-OT gradient is a reverse scan through time.  At time \(t\),
the adjoint entering the OT block contains two sources: downstream dependence
of future particles, and direct dependence of the current scalar contribution.
The OT VJP returns adjoints for the pre-transport particles and log weights:
\[
  (G_{X_t},G_{\log w_t})
  =
  \operatorname{VJP}_{\mathrm{OT},t}
  (G_{\widetilde X_t}).
\]
The log-weight adjoint then flows through the PF-PF proposal-correction terms:
transition density, observation density, proposal density, and flow
log-determinant.  The particle adjoint flows through the LEDH migration and the
model transition and observation callbacks.  Repeating this from \(t=T\) down
to \(t=1\) yields the parameter score.

\begin{proposition}[Filtering-loop custom VJP]
\label{prop:bf-ledh-ot-filtering-loop-vjp}
Suppose each time-step map in the executed LEDH-PFPF-OT value path is
differentiable under fixed branch decisions, fixed random streams, fixed masks,
and fixed iteration budgets.  If each time-step custom VJP is the VJP of that
same executed time-step map, then the reverse time scan returns the derivative
of the executed scalar \(\mathcal L_K(\theta)\).
\end{proposition}

\begin{proof}
The full filter value is a finite composition of the time-step maps and the
final scalar reduction.  By the chain rule, the VJP of the full composition is
the product of the time-step transposed Jacobians in reverse time order.  The
assumption says that each custom time-step VJP equals the corresponding
transposed Jacobian multiplication for the executed map.  Therefore their
reverse-time composition equals the VJP of the executed scalar.
\end{proof}

\section{A proposed score-aware teacher-projection extension}
\label{sec:bf-ledh-contract-e-score-aware-teacher-projection}

Contract E--Chol restores a selected value-side object: the weighted source
mean and a ridged covariance identity.  Its fixed residual design and prepared
ridge are not independent controls for likelihood accuracy, score accuracy,
and structural support.  In particular, agreement of first two moments does not
imply agreement of the tangent filtering measure.  This section therefore
studies a major algorithmic extension, not a tuning change to
\texttt{contract\_e\_chol\_v1}.

The provisional name is \emph{Contract E--teacher projection}, abbreviated
Contract E--TP.  It has no canonical, leaderboard, default, or HMC-facing status.
It changes three parts of the reset contract:
\begin{enumerate}[label=(\arabic*)]
  \item it constructs a larger finite teacher by pairing parent particles with
    fixed innovation points;
  \item it compresses selected teacher feature expectations into a smaller
    positive weighted cloud; and
  \item it carries those nonuniform weights into the next LEDH step instead of
    resetting the weights to \(1/N\).
\end{enumerate}
The current Contract E transport geometry may supply a reference cloud,
preferred anchors, or a displacement objective.  It cannot guarantee
score-feature feasibility by itself.  The existence result below instead uses
anchors selected from the teacher support.  For a structurally degenerate model
this is important: every teacher point is generated by the declared structural
map, so compression never invents an off-support full-state barycenter.

\subsection{The filtering derivative that must survive reset}

Let \(\eta_{t,\theta}\) be a normalized filtering measure.  For a fixed test
function \(h\), its tangent measure \(\zeta_{t,\theta}\) is defined weakly by
\begin{equation}
\label{eq:bf-ledh-tangent-measure-definition}
  \zeta_{t,\theta}(h)
  =D_\theta\{\eta_{t,\theta}(h)\}.
\end{equation}
The conditional path-score representation and the corresponding backward
particle recursion are given in the inspected arXiv/technical-report version of
\citet[Section~2, Eq.~(2.5), and Algorithm~1]{delmoral2015uniform}; that
algorithm evaluates the backward statistic with pairwise parent--child work.
The project construction below is different: it compresses a finite teacher
measure while preserving selected expectations and their derivatives.  The
literature supplies the filtering-derivative target; it does not prove the new
compression rule.

For a parameter-dependent feature \(h_\theta\), the total derivative is
\begin{equation}
\label{eq:bf-ledh-parameter-dependent-feature-derivative}
  D_\theta\{\eta_{t,\theta}(h_\theta)\}
  =\zeta_{t,\theta}(h_\theta)
   +\eta_{t,\theta}(D_\theta h_\theta).
\end{equation}
Here \(\zeta_{t,\theta}(h_\theta)\) means that the tangent measure at the fixed
parameter is applied to the frozen function \(x\mapsto h_\theta(x)\); the second
term accounts for the feature's explicit parameter dependence.
This distinction matters below.  Preserving a parameter-dependent predictive
potential preserves the displayed total derivative.  It does not separately
identify the two terms on the right unless the feature system also identifies
them.

\subsection{Parent-by-innovation teacher}

Suppose the incoming finite cloud has parent points and positive normalized
weights
\[
  \{(x_i,a_i)\}_{i=1}^{N_p},\qquad \sum_i a_i=1,
\]
and let a fixed innovation rule have points and positive weights
\[
  \{(\varepsilon_j,b_j)\}_{j=1}^{N_\varepsilon},\qquad
  \sum_j b_j=1.
\]
For a structural transition \(F_\theta\), form
\begin{equation}
\label{eq:bf-ledh-score-aware-teacher-candidates}
  z_{ij}(\theta)=F_\theta(x_i,\varepsilon_j).
\end{equation}
Thus a nominal \(N\)-parent by \(N\)-innovation construction has \(N^2\)
candidates.  If \(c_{t,ij}(\theta)>0\) denotes the executed current-observation
and proposal-correction factor, define
\begin{equation}
\label{eq:bf-ledh-score-aware-teacher-weights}
  u_{ij}(\theta)=a_i b_j c_{t,ij}(\theta),\qquad
  \pi_{ij}(\theta)=\frac{u_{ij}(\theta)}{\sum_{r,s}u_{rs}(\theta)}.
\end{equation}
All fixed LEDH noises, flow steps, proposal corrections, and branch choices that
belong to the finite target must be included in \(c_{t,ij}\).  Omitting them
would construct a teacher for a different scalar.

After flattening \((i,j)\) to \(a=1,\ldots,N_T\), the teacher is
\begin{equation}
\label{eq:bf-ledh-score-aware-teacher-measure}
  \eta^T_{t,\theta}=\sum_{a=1}^{N_T}\pi_a(\theta)
  \delta_{z_a(\theta)}.
\end{equation}
This is a filtering teacher, not a smoother: it uses the current filtered
parents and current innovations.  An \(O(N^2)\) backward tangent statistic may
also be evaluated on this finite system, but the score-preservation theorem
below does not inject that statistic as an externally supplied gradient.  It
differentiates the same finite primal constraints.

Choose \(K\) finite differentiable features and include mass explicitly:
\begin{equation}
\label{eq:bf-ledh-score-aware-feature-vector}
  \psi_\theta(x)
  =\left(1,x^\top,\operatorname{vech}(xx^\top)^\top,
    r_{t+1,\theta,\xi}(x),\varphi_{1,\theta}(x),\ldots\right)^\top
  \in\mathbb R^K.
\end{equation}
Here \(r_{t+1,\theta,\xi}(x)\) is the next-step finite predictive contribution
conditional on parent \(x\), with the same fixed finite noise stream \(\xi\) and
the same transition, LEDH proposal, importance correction, and observation
program that the next value evaluation will execute.  For a model transition
density it has the idealized form
\begin{equation}
\label{eq:bf-ledh-score-aware-predictive-potential}
  r_{t+1,\theta}(x)
  =\int f_\theta(x'\mid x)g_\theta(y_{t+1}\mid x')\,dx',
\end{equation}
or the exact fixed-rule finite sum corresponding to that integral.  Define the
teacher target
\begin{equation}
\label{eq:bf-ledh-score-aware-teacher-target}
  b_\theta
  =\eta^T_{t,\theta}(\psi_\theta)
  =\sum_{a=1}^{N_T}\pi_a(\theta)\psi_\theta(z_a(\theta)).
\end{equation}

\subsection{Positive compression as constrained feasibility}

Let \(s_\ell(\theta)\), \(\ell=1,\ldots,m\), be selected teacher candidates and
write
\begin{equation}
\label{eq:bf-ledh-score-aware-feature-matrix}
  \Phi_\theta
  =\begin{bmatrix}
    \psi_\theta(s_1(\theta))&\cdots&\psi_\theta(s_m(\theta))
  \end{bmatrix}.
\end{equation}
The student weights are required to satisfy
\begin{equation}
\label{eq:bf-ledh-score-aware-feasibility}
  \Phi_\theta q_\theta=b_\theta,\qquad q_\theta\geq0.
\end{equation}
Because the first feature is one, the first equation is
\(\mathbf1^\top q_\theta=1\).  Structural support, mass, moments, and predictive
features are constraints, not terms in an arbitrarily weighted multi-objective
loss.  A displacement or Contract E proximity objective may be minimized only
inside this feasible set.

\begin{proposition}[Finite positive teacher coreset]
\label{prop:bf-ledh-score-aware-positive-coreset}
For any finite teacher in
\eqref{eq:bf-ledh-score-aware-teacher-measure} and any finite \(K\)-component
feature vector whose first component is one, there exists a nonnegative student
supported on at most \(K\) teacher candidates that satisfies
\eqref{eq:bf-ledh-score-aware-feasibility} at the fixed parameter value.
\end{proposition}

\begin{proof}
The full teacher weights are already a feasible nonnegative solution because
\(b_\theta=\Phi^T_\theta\pi_\theta\).  If a feasible solution has more than \(K\)
positive entries, the corresponding columns in \(\mathbb R^K\) are linearly
dependent.  Hence there is a nonzero vector \(c\) supported on those entries
with \(\Phi_\theta c=0\).  The first row of \(\Phi_\theta\) is all ones, so
\(\mathbf1^\top c=0\); therefore \(c\) has both positive and negative entries.
Set
\[
  \alpha=\min_{c_\ell>0}\frac{q_\ell}{c_\ell},\qquad
  q'=q-\alpha c.
\]
Then \(q'\geq0\), at least one formerly positive component is zero, and
\(\Phi_\theta q'=\Phi_\theta q=b_\theta\).  Repeating the elimination produces
a feasible solution with at most \(K\) positive entries.
\end{proof}

Proposition~\ref{prop:bf-ledh-score-aware-positive-coreset} is a pointwise
existence result.  It does not say that the active support varies smoothly with
\(\theta\).  HMC cannot repeatedly choose a different linear-program basis and
then ignore those switches.  A differentiable finite target needs a fixed local
chart or a globally smooth positive projection.

\begin{proposition}[Fixed active-set differentiable chart]
\label{prop:bf-ledh-score-aware-fixed-chart}
Let \(A\) select exactly \(K\) teacher candidates.  Suppose, at a center
\(\theta_0\), that the square matrix \(\Phi_{A,\theta_0}\) is nonsingular and
\[
  q_{\theta_0}=\Phi_{A,\theta_0}^{-1}b_{\theta_0}>0
\]
componentwise.  If \(\Phi_{A,\theta}\) and \(b_\theta\) are continuously
differentiable, then in a neighborhood of \(\theta_0\) the fixed-support rule
\begin{equation}
\label{eq:bf-ledh-score-aware-square-solve}
  q_\theta=\Phi_{A,\theta}^{-1}b_\theta
\end{equation}
is positive and continuously differentiable.  For any parameter direction its
tangent is the linear solve
\begin{equation}
\label{eq:bf-ledh-score-aware-weight-tangent}
  \dot q
  =\Phi_A^{-1}\left(\dot b-\dot\Phi_A q\right).
\end{equation}
\end{proposition}

\begin{proof}
Matrix inversion is continuously differentiable on the open set of nonsingular
matrices.  Thus \(q_\theta\) is continuously differentiable near \(\theta_0\).
Strict positivity is an open condition, so it persists after reducing the
neighborhood if necessary.  Differentiating
\(\Phi_Aq=b\) yields
\(\dot\Phi_Aq+\Phi_A\dot q=\dot b\), and solving for \(\dot q\) gives
\eqref{eq:bf-ledh-score-aware-weight-tangent}.
\end{proof}

When more than \(K\) fixed anchors are desired, a semi-analytical Euclidean
projection is
\begin{equation}
\label{eq:bf-ledh-score-aware-kkt-problem}
  \min_q\ \frac12(q-q_0)^\top R(q-q_0)
  \quad\text{subject to}\quad \Phi q=b,
\end{equation}
where \(R\) is positive definite and \(q_0\) may encode Contract E transport
preferences.  If \(\Phi\) has full row rank, its equality-constrained solution is
\begin{equation}
\label{eq:bf-ledh-score-aware-kkt-solution}
  q=q_0+R^{-1}\Phi^\top
  (\Phi R^{-1}\Phi^\top)^{-1}(b-\Phi q_0).
\end{equation}
This requires a \(K\times K\) solve rather than a per-step high-dimensional
optimizer.  Positivity is not implied by
\eqref{eq:bf-ledh-score-aware-kkt-solution}; a negative component is a failed
chart, not something that may be clipped without changing the constraints.

\subsection{What the projection proves about tangents and score}

\begin{proposition}[Selected-feature tangent preservation]
\label{prop:bf-ledh-score-aware-tangent-preservation}
Under Proposition~\ref{prop:bf-ledh-score-aware-fixed-chart}, define
\[
  \eta^S_{t,\theta}=\sum_{\ell\in A}q_{\ell,\theta}
  \delta_{s_\ell(\theta)}.
\]
Then throughout the fixed-chart neighborhood,
\begin{equation}
\label{eq:bf-ledh-score-aware-feature-value-identity}
  \eta^S_{t,\theta}(\psi_\theta)
  =\eta^T_{t,\theta}(\psi_\theta),
\end{equation}
and, in every parameter direction,
\begin{equation}
\label{eq:bf-ledh-score-aware-feature-tangent-identity}
  D_\theta\{\eta^S_{t,\theta}(\psi_\theta)\}
  =D_\theta\{\eta^T_{t,\theta}(\psi_\theta)\}.
\end{equation}
For a parameter-independent retained feature \(h\), this implies
\(\zeta^S_{t,\theta}(h)=\zeta^T_{t,\theta}(h)\).
\end{proposition}

\begin{proof}
Equation~\eqref{eq:bf-ledh-score-aware-feature-value-identity} is exactly the
constraint \(\Phi_Aq=b\).  Differentiating its left side and substituting
\eqref{eq:bf-ledh-score-aware-weight-tangent} gives
\[
  \dot\Phi_Aq+\Phi_A\dot q
  =\dot\Phi_Aq+\dot b-\dot\Phi_Aq
  =\dot b,
\]
which is the derivative of the teacher target.  When \(h\) is parameter
independent, definition~\eqref{eq:bf-ledh-tangent-measure-definition} gives the
last statement.
\end{proof}

The proposition does not train a gradient to look like a teacher score.  The
same differentiable primal solve produces both the student value and its
tangent.  Consequently there is no omitted optimizer hypergradient and no
externally overwritten score.

\begin{proposition}[One-step value and score preservation]
\label{prop:bf-ledh-score-aware-next-score}
Assume the current likelihood increment \(\Delta_t(\theta)\) is computed before
reset, as required by the canonical time order.  Suppose the positive
next-step finite predictive contribution \(r_{t+1,\theta,\xi}\) is a retained
feature and is exactly the contribution executed at the next step.  Define
\[
  Z^T_{t+1}(\theta)
  =\eta^T_{t,\theta}(r_{t+1,\theta,\xi}),\qquad
  Z^S_{t+1}(\theta)
  =\eta^S_{t,\theta}(r_{t+1,\theta,\xi}).
\]
Then, throughout the fixed-chart neighborhood,
\begin{align}
  Z^S_{t+1}(\theta)&=Z^T_{t+1}(\theta),\\
  \log Z^S_{t+1}(\theta)&=\log Z^T_{t+1}(\theta),\\
  D_\theta\log Z^S_{t+1}(\theta)
    &=D_\theta\log Z^T_{t+1}(\theta).
\end{align}
Therefore the two-step partial scalars
\(\Delta_t+\log Z^S_{t+1}\) and
\(\Delta_t+\log Z^T_{t+1}\), and their scores, are identical.
\end{proposition}

\begin{proof}
The first equality is the retained predictive-feature row of
\eqref{eq:bf-ledh-score-aware-feature-value-identity}.  Positivity permits
taking logarithms.  The log values are equal as functions throughout the
fixed-chart neighborhood, so their total derivatives are equal.  The current
increment is shared because it entered the scalar before either reset.
\end{proof}

This result is exact relative to the declared finite teacher and selected
one-step feature.  It is not a proof of equality to the exact nonlinear filter,
nor does it determine all future likelihood functionals.  A full recursive
method must construct and validate a new teacher and feature chart at every
reset while retaining a globally differentiable finite target.  Whether a
single fixed chart remains positive over an HMC-relevant parameter region is an
open admission question.

\subsection{Progressive score information and compact score-aware features}
\label{sec:bf-ledh-tp-progressive-score-features}

The one-step proposition above explains both the success of a two-observation
test and its limitation.  The feature \(r_{t+1,\theta}\) identifies one future
normalizing constant, but it does not identify how score information accumulated
before time \(t\) is distributed over the current state.  A scalar cumulative
score cannot repair this loss: it is constant across particles and therefore
adds no constraint beyond mass.  The state-dependent object is the conditional
additive score statistic.

For clarity, first consider a hidden Markov model with initial density
\(\mu_\theta\), transition density \(f_\theta\), and observation density
\(g_\theta\).  With observations through time \(t\), define the complete-data
additive score
\begin{equation}
\label{eq:bf-ledh-progressive-complete-score}
  H_{t,\theta}(x_{0:t})
  =\nabla_\theta\log\mu_\theta(x_0)
   +\sum_{s=1}^{t}\nabla_\theta\log f_\theta(x_s\mid x_{s-1})
   +\sum_{s=0}^{t}\nabla_\theta\log g_\theta(y_s\mid x_s),
\end{equation}
and its state-conditioned version
\begin{equation}
\label{eq:bf-ledh-progressive-score-mark}
  \tau_{t,\theta}(x)
  =\mathbb E_\theta\!\left[
      H_{t,\theta}(X_{0:t})\mid X_t=x,y_{0:t}
    \right].
\end{equation}
This is a \(p\)-vector when \(\theta\in\mathbb R^p\).  Its filtering mean is the
observed-data score
\(S_{t,\theta}=\eta_{t,\theta}(\tau_{t,\theta})\).  The centered progressive
score field is
\begin{equation}
\label{eq:bf-ledh-centered-progressive-score}
  c_{t,\theta}(x)=\tau_{t,\theta}(x)-S_{t,\theta},
  \qquad \eta_{t,\theta}(c_{t,\theta})=0.
\end{equation}
This conditional expectation is the finite-dimensional statistic used in the
backward particle interpretation of the filter derivative
\citep[Section~2, Eq.~(2.5), Algorithm~1]{delmoral2015uniform}.  Computing it
does not require sampling complete trajectories or storing their histories, but
the exact finite-particle update uses pairwise parent--child work.

\begin{proposition}[Tangent measure represented by the centered score field]
\label{prop:bf-ledh-progressive-score-tangent-representation}
Assume differentiation may be interchanged with the state integrals and all
displayed expectations are finite.  For every parameter-independent integrable
test function \(h\),
\begin{equation}
\label{eq:bf-ledh-progressive-score-covariance}
  \zeta_{t,\theta}(h)
  =\nabla_\theta\{\eta_{t,\theta}(h)\}
  =\eta_{t,\theta}\!\left(hc_{t,\theta}^{\top}\right).
\end{equation}
For a differentiable parameter-dependent \(h_\theta\),
\begin{equation}
\label{eq:bf-ledh-progressive-score-total-derivative}
  \nabla_\theta\{\eta_{t,\theta}(h_\theta)\}
  =\eta_{t,\theta}(\nabla_\theta h_\theta)
   +\eta_{t,\theta}\!\left(h_\theta c_{t,\theta}^{\top}\right).
\end{equation}
\end{proposition}

\begin{proof}
Write the normalized filtering density as
\(\eta_{t,\theta}(x_t)=\gamma_{t,\theta}(x_t)/Z_{t,\theta}\).
Differentiating the unnormalized path integral for \(\gamma_{t,\theta}\)
multiplies its integrand by the complete-data score
\eqref{eq:bf-ledh-progressive-complete-score}.  Conditioning on \(X_t=x_t\)
therefore gives
\(\nabla_\theta\log\gamma_{t,\theta}(x_t)
  =\tau_{t,\theta}(x_t)\).
Integrating this identity gives
\(\nabla_\theta\log Z_{t,\theta}=S_{t,\theta}\), hence
\(\nabla_\theta\eta_{t,\theta}(x)
  =\eta_{t,\theta}(x)c_{t,\theta}(x)\).
Integration against \(h\) proves
\eqref{eq:bf-ledh-progressive-score-covariance}.  The product rule adds
\(\eta_{t,\theta}(\nabla_\theta h_\theta)\) when the feature depends explicitly
on the parameter, proving
\eqref{eq:bf-ledh-progressive-score-total-derivative}.
\end{proof}

The discrete \(O(N^2)\) recursion is explicit for an ordinary particle
approximation of the model filter.  Given parents
\((x_i,a_i,\tau_i)_{i=1}^{N_p}\) and common child locations \(z_k\), define the
normalized backward weights
\begin{equation}
\label{eq:bf-ledh-progressive-backward-weight}
  B_{ki}
  =\frac{a_i f_\theta(z_k\mid x_i)}
         {\sum_{j=1}^{N_p}a_j f_\theta(z_k\mid x_j)}.
\end{equation}
Then the child mark is
\begin{equation}
\label{eq:bf-ledh-progressive-score-recursion}
  \tau'_k
  =\nabla_\theta\log g_\theta(y_{t+1}\mid z_k)
   +\sum_{i=1}^{N_p}B_{ki}
      \left\{\tau_i+
        \nabla_\theta\log f_\theta(z_k\mid x_i)\right\}.
\end{equation}
Initial-law score terms enter the initial marks.

Equations~\eqref{eq:bf-ledh-progressive-backward-weight}--%
\eqref{eq:bf-ledh-progressive-score-recursion} must not be copied unchanged
into the corrected LEDH finite program.  Its teacher candidates are
parent-specific deterministic quadrature/proposal images, so their discrete
backward kernel is not automatically the continuous model kernel \(f_\theta\).
A finite-target mark must be derived from the exact executed proposal measure,
including transition, observation, proposal correction, Jacobian, quadrature
mass, and support-motion terms.  Using only \(f_\theta\), or treating the
derivative of a moving atomic measure as an ordinary density score, computes a
different object.  Until that derivation is supplied, model-filter score fields
may be used only as labeled diagnostic features.

The full tensor-product proposal
\(\{\phi_j(x)c_k(x):j=1,\ldots,K_0,\ k=1,\ldots,p\}\) uses \(K_0p\)
additional constraints and is generally too large.  A compact diagnostic
candidate retains
\begin{equation}
\label{eq:bf-ledh-compact-progressive-feature-vector}
  \Psi_{t,\theta}(x)
  =\left(
    \phi_{t,\theta}(x)^{\top},
    c_{t,\theta}(x)^{\top},
    \{r_{t+1,\theta}(x)c_{t,\theta}(x)\}^{\top}
   \right)^{\top}.
\end{equation}
The displayed form adds \(2p\), rather than \(K_0p\), rows.  There is a smaller
gauge-fixed implementation.  Replacing every parent mark by
\(\tau_i+a\), for a common vector \(a\), replaces every child mark in
\eqref{eq:bf-ledh-progressive-score-recursion} by \(\tau'_k+a\), because each
backward row sums to one.  Subsequent centering removes \(a\) exactly.  Hence an
implementation that carries centered marks and explicitly re-centers them after
selection need not impose the \(p\) zero-mean rows through the projection solve;
it verifies zero mean as an invariant instead.  The nonredundant compact chart
then appends only \(r_{t+1,\theta}c_{t,\theta}\), using \(K_0+p\) rows.  This asks
the coreset to retain the score-bearing interaction with the next potential.  It
is a feature-sufficiency hypothesis, not yet an exact finite-LEDH tangent
theorem.

At the first LGSSM reset this reduction is also forced by rank: the observation
scale component of \(c_{t,\theta}\) is a quadratic function of the current state
and lies in the mass/first/second-moment span.  Retaining all \(2p\) rows makes
that feature matrix singular.  Dropping the zero-mean rows by the gauge argument
avoids data-dependent feature deletion and leaves a fixed 16-row LGSSM chart.

\begin{proposition}[One-step historical-score identity for a dominated filter]
\label{prop:bf-ledh-progressive-next-score-preservation}
Let teacher and student be two dominated filtering measures whose tangent
measures admit Proposition~\ref{prop:bf-ledh-progressive-score-tangent-representation}
with the same declared score-coordinate convention.  Suppose
\begin{align}
  \eta_t^S(r_{t+1,\theta})&=\eta_t^T(r_{t+1,\theta}),
    \label{eq:bf-ledh-progressive-r-match}\\
  \eta_t^S(r_{t+1,\theta}c_S^{\top})
    &=\eta_t^T(r_{t+1,\theta}c_T^{\top}),
    \label{eq:bf-ledh-progressive-rc-match}\\
  \eta_t^S(\nabla_\theta r_{t+1,\theta})
    &=\eta_t^T(\nabla_\theta r_{t+1,\theta}).
    \label{eq:bf-ledh-progressive-dr-match}
\end{align}
Then the next normalizing constant and its total score agree.
\end{proposition}

\begin{proof}
The first identity gives the common positive normalizing constant.  Applying
\eqref{eq:bf-ledh-progressive-score-total-derivative} to
\(h_\theta=r_{t+1,\theta}\) shows that its derivative is the sum of the
explicit-feature term in \eqref{eq:bf-ledh-progressive-dr-match} and the
historical tangent term in \eqref{eq:bf-ledh-progressive-rc-match}.  Both agree,
so division by the common normalizing constant proves equality of the
log-score.
\end{proof}

The proposition applies to the dominated model-filter setting just stated.  A
moving discrete Contract E--TP cloud must first prove its corresponding tangent
representation; matching sampled \(c\) and \(rc\) rows alone does not prove that
its autodiff tangent measure equals \(q\,c\).  Consequently the compact rows can
diagnose and potentially reduce information loss, but their success must still
be judged by the total derivative of the executed scalar.

For an LGSSM there is a stronger finite-horizon diagnostic.  The remaining-data
likelihood
\begin{equation}
\label{eq:bf-ledh-lgssm-continuation-feature}
  R_{t+1:T,\theta}(x)
  =p_\theta(y_{t+1:T}\mid X_t=x)
\end{equation}
is available from a conditional Kalman recursion and is an exponential
quadratic function of \(x\).  It is one parameter-dependent scalar feature
regardless of \(p\) and \(T\).
Writing its stabilized ratio as
\[
 \frac{R_{t+1:T,\theta}(x)}{R_{t+1:T,\theta}(0)}
 =\exp\{-\tfrac12x^\top J_{t,\theta}x+h_{t,\theta}^\top x\},
\]
all pairs \((J_{t,\theta},h_{t,\theta})\) are obtained in one reverse-time
Gaussian information recursion.  This avoids rebuilding a length-\(T-t\)
conditional Kalman graph at every reset: construction is \(O(Td_x^3)\), followed
by \(O(Nd_x^2)\) evaluation per cloud.  The stabilization removes only a common
positive factor from the feature row and therefore does not change its linear
constraint.

\begin{proposition}[Finite-horizon continuation identity]
\label{prop:bf-ledh-lgssm-continuation-certificate}
Fix \(T\) and a parameter neighborhood on which one positive nonsingular chart
is used.  If, immediately after reset at time \(t\),
\begin{equation}
  \eta^S_{t,\theta}(R_{t+1:T,\theta})
  =\eta^T_{t,\theta}(R_{t+1:T,\theta})
\end{equation}
as an identity throughout that neighborhood, then teacher and student assign
the same exact-model remaining-data likelihood and total derivative to their
respective time-\(t\) measures.
\end{proposition}

\begin{proof}
For either measure, the displayed expectation is the exact-model conditional
factorization of the future observations given its time-\(t\) state
distribution.  The two sides are equal positive differentiable functions of
\(\theta\) on the fixed-chart neighborhood.  Taking logarithms and
differentiating the same identity proves equality of their total derivatives,
including both measure dependence and explicit parameter dependence of
\(R_{t+1:T,\theta}\).
\end{proof}

This does not prove equality of the later recursively executed finite LEDH
scalar.  Its quadrature/proposal operator differs from exact Kalman integration,
and subsequent projections are nonlinear cloud maps.  The continuation row is
therefore an LGSSM oracle diagnostic of reset information loss.  It may test
whether a compact future-information summary repairs the observed drift, but it
cannot be transferred to nonlinear models without a separately justified
finite continuation approximation.  Neither this row nor the model-filter
score field replaces the missing finite-LEDH tangent derivation.

\subsubsection{A bounded ladder of progressive continuation marks}
\label{sec:bf-ledh-progressive-continuation-ladder}

One full remaining-data continuation feature is compact in row count, but its
offline construction can be expensive and its approximation can be fragile in
a persistent nonlinear model.  Replacing it by the entire sequence of all
future prefixes would make the chart grow with the horizon.  The implemented
intermediate design instead freezes a small ordered set of positive look-ahead
lengths
\(
  \mathcal L=\{\ell_1<\cdots<\ell_r\}
\)
before preparation.  For each available prefix define
\begin{equation}
\label{eq:bf-ledh-progressive-continuation-mark}
  B_{t,\ell,\theta}(x)
  =\int
    \prod_{s=t+1}^{t+\ell}
      f_\theta(x_s\mid x_{s-1})
      g_\theta(y_s\mid x_s)
    \,d x_{t+1:t+\ell},
  \qquad x_t=x.
\end{equation}
For a finite record ending at \(T\), the realized length is
\(
  \bar\ell_t=\min(\ell,T-t)
\).
If two requested lengths have the same realized length, only one row is kept.
This de-duplication depends only on \(t,T,\mathcal L\), not on particles,
parameters, observations, or numerical values, so it is a prepared chart shape
rather than a runtime active-set decision.

Let \(x_\star(\theta)\) be the declared reference point and set
\begin{equation}
\label{eq:bf-ledh-progressive-continuation-stabilized-mark}
  h_{t,\ell,\theta}(x)
  =\frac{B_{t,\ell,\theta}(x)}
         {B_{t,\ell,\theta}(x_\star(\theta))}.
\end{equation}
The implementation evaluates numerator and denominator by the same fixed-grid
backward target recursion.  The denominator is positive and common to a row;
it changes neither the row span nor the equality constraint.  Its dependence
on \(\theta\), including dependence through \(x_\star(\theta)\), remains in the
autodiff graph.  For scalar SV the feature vector at time \(t\) is
\begin{equation}
\label{eq:bf-ledh-progressive-sv-feature-vector}
  \psi_{t,\theta}(x)
  =\left(1,x,x^2,
    \{h_{t,\bar\ell,\theta}(x):
       \bar\ell\in\overline{\mathcal L}_t\}\right)^\top,
\end{equation}
where \(\overline{\mathcal L}_t\) contains the distinct realized lengths.
Thus the reset carries at most \(3+r\) weights and scalar locations, independent
of \(T\); it does not carry a growing score history.  The backward preparation
cost still grows with each requested continuation window, which is why
\(\mathcal L\) is a capacity ladder rather than an unrestricted feature list.

\subsubsection{Actual-SV amendment: an overcomplete analytical chart}
\label{sec:bf-ledh-actual-sv-overcomplete-chart}

A square chart uses one anchor for each feature.  It is algebraically unique
when its feature matrix is nonsingular, but its positive domain can be much
smaller than the parameter neighborhood on which a score is needed.  This is
not repaired by adding features: each new feature adds an equality constraint.
The amendment here keeps the feature target fixed and adds anchors, thereby
adding degrees of freedom to its representation.

The failure that motivates the amendment is concrete.  The Actual-SV
experiment uses the two unconstrained coordinates
\[
  \theta=(\Phi^{-1}(\gamma),\log\beta),\qquad
  \theta_0=(0.2533471031357997,-0.916290731874155),
\]
and, at every nonterminal time, the four features
\begin{equation}
\label{eq:bf-ledh-actual-sv-four-features}
  \psi_{t,\theta}(x)
  =\bigl(1,x,x^2,h_{t,16,\theta}(x)\bigr)^\top,
  \qquad q=4.
\end{equation}
The center-prepared four-anchor chart at zero-based time index \(748\) had
minimum weight \(3.58612667\times10^{-7}\).  Holding its prepared indices and
row scales fixed, changing the first coordinate by \(-10^{-5}\) produced
\(-4.22267303\times10^{-6}\), while changing the second coordinate by
\(+10^{-5}\) produced \(-5.23620186\times10^{-6}\).  Thus the executed square
chart correctly failed closed: those moment targets are outside the positive
simplex generated by those four frozen columns.  This is a chart failure, not
a failure of the SV parameter transform, Sinkhorn balancing, or canonical
Contract E--Chol.

Choose one anchor count \(K>q\), constant across all nonterminal times.  At a
given time, preparation freezes teacher indices \(i_1,\ldots,i_K\).  The
indices are discrete prepared data, but the selected locations
\(z_k(\theta)=z_{i_k}(\theta)\) are not frozen: they move with the finite
teacher program.  Consequently their feature columns contribute to the total
derivative.  Let \(D\in\R^{q\times q}\) be a positive diagonal row-scale
matrix prepared at the center and then frozen, and define
\begin{equation}
\label{eq:bf-ledh-actual-sv-scaled-system}
  A_\theta=D^{-1}
  \begin{bmatrix}
    \psi_{t,\theta}(z_1(\theta))&\cdots&
    \psi_{t,\theta}(z_K(\theta))
  \end{bmatrix},
  \qquad
  b_\theta=D^{-1}m_{t,\theta},
\end{equation}
where \(m_{t,\theta}\) is the same finite teacher feature target used by the
square chart.  Positivity and nonsingularity of \(D\) make row scaling an
invertible numerical transformation, not a change in the constraints.

\paragraph{Center reference.}
The anchor indices are deterministic center-teacher weighted quantiles.  For
anchor count \(K\), their probability levels are
\((k-\tfrac12)/K\), \(k=1,\ldots,K\).  Stable sorting, cumulative teacher mass,
and left search select indices.  Repeated indices are removed stably and, only
to fill duplicates, unused teacher indices are added in descending normalized
mass order with stable tie-breaking.  Every center teacher point is then
assigned to its nearest anchor in state space, again with stable tie-breaking.
Aggregating its normalized mass gives a vector \(v\in\R^K\).  Preparation
rejects an empty cell: \(v_k>0\) for every \(k\) is a hard condition and is
never manufactured by clipping or flooring.

Write \(V=\diag(v)\), \(A_0=A_{\theta_0}\), and \(b_0=b_{\theta_0}\).  The
center reference is the Pearson projection
\begin{equation}
\label{eq:bf-ledh-actual-sv-reference-projection}
  r=v+VA_0^\top(A_0VA_0^\top)^{-1}(b_0-A_0v).
\end{equation}

\begin{proposition}[Analytical Pearson center reference]
\label{prop:bf-ledh-actual-sv-reference-projection}
Suppose \(v>0\) componentwise and \(A_0\in\R^{q\times K}\) has full row rank.
Then \eqref{eq:bf-ledh-actual-sv-reference-projection} is the unique solution of
\begin{equation}
\label{eq:bf-ledh-actual-sv-reference-program}
  \min_{u\in\R^K}
  \frac12(u-v)^\top V^{-1}(u-v)
  \quad\text{subject to}\quad A_0u=b_0.
\end{equation}
It satisfies \(A_0r=b_0\).  Strict positivity of \(r\) does not follow from the
equality problem and is therefore a separate preparation gate.
\end{proposition}

\begin{proof}
The Hessian \(V^{-1}\) is positive definite, so the objective is strictly
convex.  The Lagrange equations for multiplier \(\nu\) are
\[
  V^{-1}(u-v)+A_0^\top\nu=0,
  \qquad A_0u=b_0.
\]
Eliminating \(u=v-VA_0^\top\nu\) gives
\(
  (A_0VA_0^\top)\nu=A_0v-b_0
\).
Full row rank of \(A_0\) and \(V\succ0\) make this Gram matrix positive
definite.  Substitution yields
\eqref{eq:bf-ledh-actual-sv-reference-projection}; it is feasible and, by
strict convexity, unique.  Equality feasibility alone imposes no sign
constraint, so a nonpositive component must be rejected rather than hidden.
\end{proof}

For a positive accepted \(r\), freeze the Pearson precision and its inverse,
\begin{equation}
\label{eq:bf-ledh-actual-sv-pearson-metric}
  P=\diag(r_1^{-1},\ldots,r_K^{-1}),
  \qquad R=P^{-1}=\diag(r).
\end{equation}
The choice is the Hessian of the Pearson chi-square discrepancy centered at
\(r\); it is a declared metric hypothesis, not a statistical confidence
threshold.  Runtime computes
\begin{equation}
\label{eq:bf-ledh-progressive-kkt-program}
  w_\theta
  =\arg\min_w\frac12(w-r)^\top P(w-r)
  \quad\text{subject to}\quad A_\theta w=b_\theta.
\end{equation}

\begin{proposition}[Runtime solution and frozen-reference total JVP]
\label{prop:bf-ledh-progressive-kkt-chart}
Suppose \(r>0\), \(A_\theta\) has full row rank, and \(A_\theta,b_\theta\) are
differentiable along a parameter direction.  Set
\begin{equation}
\label{eq:bf-ledh-progressive-kkt-solution}
  S=A_\theta R A_\theta^\top,\qquad
  c=b_\theta-A_\theta r,\qquad
  \lambda=S^{-1}c,\qquad
  w_\theta=r+RA_\theta^\top\lambda.
\end{equation}
Then \(w_\theta\) is the unique minimizer of
\eqref{eq:bf-ledh-progressive-kkt-program} and satisfies
\(A_\theta w_\theta=b_\theta\).  With the preparation
\(r,R,D,i_1,\ldots,i_K\) frozen, its directional derivative is
\begin{align}
  \dot S&=\dot ARA^\top+AR\dot A^\top,&
  \dot c&=\dot b-\dot A r,
  \label{eq:bf-ledh-actual-sv-kkt-jvp-a}\\
  \dot\lambda&=S^{-1}(\dot c-\dot S\lambda),&
  \dot w&=R(\dot A^\top\lambda+A^\top\dot\lambda).
  \label{eq:bf-ledh-actual-sv-kkt-jvp-b}
\end{align}
In particular,
\begin{equation}
\label{eq:bf-ledh-actual-sv-differentiated-constraint}
  A\dot w+\dot A w=\dot b,
\end{equation}
so the same finite program matches both the declared features and their total
directional derivatives.
\end{proposition}

\begin{proof}
The Lagrange equations for
\eqref{eq:bf-ledh-progressive-kkt-program}, using the sign convention
\(P(w-r)-A^\top\lambda=0\), give
\(w=r+RA^\top\lambda\).  Substitution into \(Aw=b\) gives
\(S\lambda=c\).  Since \(R\succ0\) and \(A\) has full row rank, \(S\succ0\),
so the multiplier and the strictly convex minimizer are unique.  This proves
\eqref{eq:bf-ledh-progressive-kkt-solution} and exact feasibility.

Differentiate \(S\lambda=c\) while \(r\) and \(R\) are fixed:
\(
  S\dot\lambda+\dot S\lambda=\dot c
\).
Solving for \(\dot\lambda\), and then differentiating
\(w=r+RA^\top\lambda\), yields
\eqref{eq:bf-ledh-actual-sv-kkt-jvp-a}--
\eqref{eq:bf-ledh-actual-sv-kkt-jvp-b}.  Finally, differentiating the already
proved identity \(Aw=b\) gives
\eqref{eq:bf-ledh-actual-sv-differentiated-constraint}.  Here \(\dot A\)
contains the explicit parameter derivative of each feature, teacher motion,
and motion of every selected location \(z_{i_k}(\theta)\); freezing an index
does not freeze the tensor gathered at that index.
\end{proof}

\paragraph{Finite-program gate.}
The claimed object is the executed TensorFlow solve, not an unknown exact-real
solution.  Every prepared and runtime case must have finite inputs and outputs,
full row rank, condition--roundoff below one, scaled equality residual no
larger than its forward-error bound, and strictly positive computed weights.
The implementation solves the \(q\times q\) Gram system by Cholesky or a linear
solve; it never forms a matrix inverse and never stores or solves a dense
\(K\times K\) precision system.  Failed cases return an invalid flag and poison
claim-bearing outputs.  No clipping, weight floor, feature deletion, runtime
optimization, or active-set switch is allowed.

These diagnostics identify an unreliable finite solve but are not a rigorous
componentwise sign certificate for the exact-real solution.  For a selected
capacity, the weakest design case and weakest held-out case---defined by the
smallest finite TensorFlow weight, then time and predeclared point order---are
also recomputed at increasing CPU precision.  Stable sign is corroborating
evidence only, not interval arithmetic and not part of the BayesFilter
algorithmic backend.

\begin{algorithm}[H]
\caption{Actual-SV overcomplete analytical chart preparation and runtime}
\label{alg:bf-ledh-actual-sv-overcomplete-chart}
\begin{algorithmic}[1]
  \State \textbf{Prepare:} freeze one global \(K>q\), the center
    \(\theta_0\), row scales \(D_t\), and the design and held-out parameter
    points before observing candidate results.
  \For{each nonterminal time \(t\)}
    \State Select \(K\) center-teacher weighted-quantile indices with the
      deterministic duplicate-fill rule.
    \State Aggregate stable nearest-anchor Voronoi masses \(v_t\); reject if
      any component is nonfinite or nonpositive.
    \State Form \(A_{t,0},b_{t,0}\) and solve
      \eqref{eq:bf-ledh-actual-sv-reference-projection}; reject unless rank,
      conditioning, residual, finiteness, and \(r_t>0\) gates pass.
    \State Freeze indices, \(D_t\), \(r_t\), and \(R_t=\diag(r_t)\).
  \EndFor
  \State \textbf{Run:} gather the moving locations and feature columns at the
    frozen indices, form \(A_{t,\theta},b_{t,\theta}\), and solve the
    \(q\times q\) system in
    \eqref{eq:bf-ledh-progressive-kkt-solution}.
  \State Apply every finite-program gate; carry the resulting nonuniform
    positive weights, or poison the claim-bearing recursion on failure.
  \State For each parameter direction, propagate
    \eqref{eq:bf-ledh-actual-sv-kkt-jvp-a}--
    \eqref{eq:bf-ledh-actual-sv-kkt-jvp-b} through the same fixed-shape
    filtering loop.
\end{algorithmic}
\end{algorithm}

For \(q\) features and \(K\) anchors, forming \(ARA^\top\) costs
\(O(Kq^2)\), its solve costs \(O(q^3)\), and temporary storage is
\(O(Kq+q^2)\).  In Actual SV, \(q=4\), so increasing anchor capacity does not
increase the dimension of the linear solve.  The recursive GPU route must use
a fixed-shape \texttt{tf.while\_loop}; Python horizon unrolling remains a
reference diagnostic only.

This amendment proves the analytical projection, equality identities, and JVP
of the declared finite chart.  It does not prove that a tested finite set of
parameter points covers an HMC trajectory, that the four-feature teacher is an
exact nonlinear filter, that the resulting score agrees with an independent
dense target, or that Contract E--TP is canonical, leaderboard-ready, or
superior.  Those are separate empirical or scientific questions.  In
particular, matching one continuation mark and its derivative does not identify
the entire filtering tangent measure; recursive same-target and quadrature
refinement comparisons remain necessary descriptive evidence.

\subsubsection{Executed Actual-SV selection and certification record}
\label{sec:bf-ledh-actual-sv-overcomplete-certification}

The preceding propositions state what follows mathematically from an accepted
fixed preparation.  This subsection records separately how the preparation was
selected, which finite cases were executed, which derivative references were
usable, and what the resulting evidence can and cannot establish.  Numerical
observations below are not assumptions of
Propositions~\ref{prop:bf-ledh-actual-sv-reference-projection} and
\ref{prop:bf-ledh-progressive-kkt-chart}.

\paragraph{Frozen parameter sets and capacity selection.}
Write
\[
  \theta(z)=\theta_0+10^{-5}z,
  \qquad z\in\R^2.
\]
Before the capacity ladder was evaluated, the ordered design set was frozen as
\begin{equation}
\label{eq:bf-ledh-actual-sv-design-set}
 \mathcal Z_{\rm design}
 =\{(0,0),(-1,0),(1,0),(0,-1),(0,1),
      (-1,-1),(-1,1),(1,-1),(1,1)\},
\end{equation}
and the disjoint audit set was frozen as
\begin{equation}
\label{eq:bf-ledh-actual-sv-heldout-set}
 \mathcal Z_{\rm audit}
 =\{(-\tfrac12,-\tfrac12),(-\tfrac12,\tfrac12),
      (\tfrac12,-\tfrac12),(\tfrac12,\tfrac12),
      (-\tfrac34,-\tfrac14),(-\tfrac34,\tfrac14),
      (\tfrac34,-\tfrac14),(\tfrac34,\tfrac14)\}.
\end{equation}
The audit set did not participate in capacity selection and could not authorize
trying another capacity after a failure.  One global capacity was required at
all \(999\) nonterminal times.  The selection rule was the smallest
\(K\in\{5,\ldots,25\}\) passing every design point at each of
\(T=10,100,1000\).

\begin{table}[H]
\centering
\caption{Frozen Actual-SV anchor-capacity ladder.  A failure is classified by
the first mathematical preparation gate it violates, not by timeout.}
\label{tab:bf-ledh-actual-sv-capacity-ladder}
\resizebox{\textwidth}{!}{%
\begin{tabular}{@{}lll@{}}
\toprule
Capacity & Observed disposition & Classification \\
\midrule
\(K=5,6\) & center preparation fails by \(T=10\) & negative Pearson reference \\
\(K=7,\ldots,22\) & center preparation passes \(T=10\), fails by \(T=100\) & negative Pearson reference \\
\(K=23\) & pass all design points at \(T=10,100,1000\) & selected smallest pass \\
\(K=24,25\) & fail at time zero & zero Voronoi/reference mass \\
\bottomrule
\end{tabular}
}
\end{table}

The selected \(K=23\) preparation passed all eight audit points once, without
retuning.  Its weakest design weight was
\(5.395467979032687\times10^{-165}\), at zero-based ordered design point \(7\), time zero,
anchor \(22\).  Its weakest audit weight was
\(5.396877738726166\times10^{-165}\), at
zero-based ordered audit point \(2\), \(z=(\tfrac12,-\tfrac12)\), time zero,
anchor \(22\).  Re-solving the saved
\(4\)-by-\(23\) systems with \(50\), \(100\), and \(200\) decimal digits kept
the same weakest anchor and a strictly positive sign.  The \(200\)-digit design
value begins \(5.39546797902344038197\times10^{-165}\), and the corresponding
audit value begins \(5.39687773872662767248\times10^{-165}\).  This supports
the sign of the saved cases; it is not interval arithmetic and does not certify
nearby unsampled parameters.

\begin{proposition}[Scope of a finite chart audit]
\label{prop:bf-ledh-actual-sv-finite-audit-scope}
Let \(\mathcal Z\) be a finite parameter set.  If the fixed prepared chart
passes every chart finite-program gate at every \(z\in\mathcal Z\), then the
executed chart and its claim-bearing scalar outputs are admissible at those
tested points.  At a tested point where the separate own-scalar derivative
gate also passes, the checked JVP is admissible for that same finite scalar.
These observations alone do not imply that either class of gate passes on
\(\operatorname{conv}(\mathcal Z)\), on any open neighborhood, or along an HMC
trajectory.
\end{proposition}

\begin{proof}
The first two statements are exactly the corresponding conjunctions of the
pointwise chart and derivative gate outcomes.
A finite set contains no information about an untested point without an
additional analytic bound, interval certificate, or verified continuity
margin.  A continuous weight can be positive at every tested point and
negative between two of them.  Therefore the stronger region claims do not
follow from the finite audit.
\end{proof}

\paragraph{Own-scalar derivative references.}
For the two-parameter center score, define the central finite-difference vector
with \(h=10^{-5}\) by
\begin{equation}
\label{eq:bf-ledh-actual-sv-fd-score}
  g^{\rm FD}_j
  =\frac{\mathcal L(\theta_0+h e_j)
          -\mathcal L(\theta_0-h e_j)}{2h},
  \qquad j=1,2,
\end{equation}
and report
\begin{equation}
\label{eq:bf-ledh-actual-sv-fd-relative-error}
  E_{\rm FD}
  =\frac{\lVert g^{\rm manual}-g^{\rm FD}\rVert_2}
         {\max(\lVert g^{\rm FD}\rVert_2,1)}.
\end{equation}
The frozen FD-only diagnostic tolerance was
\(0.05\sqrt p=0.05\sqrt2=0.07071067811865477\).  This is an owner-selected
engineering tolerance for this deterministic directional FD comparison.  It
is \emph{not} a \(95\%\) confidence interval, a two-standard-deviation rule, or
a statistical error probability: neither side of
\eqref{eq:bf-ledh-actual-sv-fd-relative-error} is a sample from a declared
sampling distribution.

\begin{table}[H]
\centering
\caption{Actual-SV manual total JVP against central finite differences of the
same finite scalar.}
\label{tab:bf-ledh-actual-sv-fd-results}
\begin{tabular}{@{}rll@{}}
\toprule
\(T\) & \(E_{\rm FD}\) & FD-only gate \\
\midrule
2    & \(2.323339358376968\times10^{-11}\) & pass \\
10   & \(1.6187801132686436\times10^{-10}\) & pass \\
100  & \(2.247653507298375\times10^{-9}\) & pass \\
1000 & \(5.7988816510053705\times10^{-8}\) & pass \\
\bottomrule
\end{tabular}
\end{table}

Ordinary TensorFlow reverse-mode autodiff agrees at the center through
\(T=10\), but one tested \(T=10\) endpoint is already nonfinite and all five
tested points are nonfinite at \(T\geq100\), even though the primal, manual JVP,
and FD values remain finite.  Component-level investigations associate the
failure with reverse propagation through extremely small positive-weight
paths; they do not prove a unique failing operation.  The observed nonfinite
result is therefore an implementation failure of that reverse-mode oracle for
this finite program.  By itself it is evidence neither that the finite scalar
is nondifferentiable nor that the manual JVP is correct.  A trial custom reverse
pullback also returned nonfinite results at \(T=100\); it was removed rather
than retained as a fallback.

An independent TensorFlow eager \texttt{ForwardAccumulator} differentiated the
unmodified primal instead.  With

\[
  E_{\rm FAD}
  =\frac{\lVert g^{\rm manual}-g^{\rm FAD}\rVert_2}
         {\max(\lVert g^{\rm manual}\rVert_2,
                    \lVert g^{\rm FAD}\rVert_2,10^{-12})},
\]
the observed comparisons were:

\begin{table}[H]
\centering
\caption{Actual-SV manual total JVP against eager forward autodiff of the same
unmodified scalar.  The reference gate is
\(\sqrt{\epsilon_{64}}=1.4901161193847656\times10^{-8}\).}
\label{tab:bf-ledh-actual-sv-forward-ad-results}
\begin{tabular}{@{}rll@{}}
\toprule
\(T\) & \(E_{\rm FAD}\) & float64 reference gate \\
\midrule
10   & \(1.298650134071988\times10^{-13}\) & pass \\
100  & \(4.447444400460831\times10^{-12}\) & pass \\
1000 & \(2.6364004959309425\times10^{-12}\) & pass \\
\bottomrule
\end{tabular}
\end{table}

At \(T=1000\), the manual and eager-forward scores were, respectively,
\begin{align*}
 g^{\rm manual}&=(5.664478697500771,-2.566049221886212),\\
 g^{\rm FAD}&=(5.664478697506889,-2.5660492218710016).
\end{align*}
Their maximum absolute difference was
\(1.5210499526574495\times10^{-11}\).  Automatic forward mode is a CPU
reference exception, not the production route.  XLA compilation of the
transformed nested loop failed on a nonconstant transpose permutation; a
non-JIT \texttt{tf.function} transformation instead failed a TensorList shape
invariant.  The production score is consequently the explicit two-direction
manual JVP inside the XLA-default fixed-shape loop.  This choice is supported
by the reported same-scalar checks; it is not justified merely by the failure
of the other autodiff implementations.

\paragraph{Independent dense-filter comparison.}
Own-scalar differentiation does not establish that the scalar approximates the
desired nonlinear filtering value or score.  The selected candidate was
therefore compared at the center with a deterministic dense Actual-SV filter
using Legendre orders \(129\) and \(257\) on radius \(10\).  For scalar score
components \(a,b\), the reported symmetric-relative difference is
\begin{equation}
\label{eq:bf-ledh-actual-sv-symmetric-relative}
  d_{\rm sym}(a,b)
  =\begin{cases}
     2|a-b|/(|a|+|b|),& |a|+|b|>0,\\
     0,&a=b=0.
   \end{cases}
\end{equation}

\begin{table}[H]
\centering
\caption{Descriptive candidate-minus-order-257 dense-reference comparisons.
No scientific-equivalence threshold was declared.}
\label{tab:bf-ledh-actual-sv-dense-comparison}
\resizebox{\textwidth}{!}{%
\begin{tabular}{@{}rlll@{}}
\toprule
\(T\) & Value difference & Absolute score difference & Componentwise \(d_{\rm sym}\) \\
\midrule
2 & \(1.3014182620452175\times10^{-8}\)
  & \((2.6158827104438842\times10^{-8},\,2.8765310133849198\times10^{-9})\)
  & \((1.187693931512109\times10^{-7},\,5.435712744272272\times10^{-9})\) \\
10 & \(2.4500970141616563\times10^{-8}\)
  & \((1.604391193232857\times10^{-7},\,7.053514550214857\times10^{-9})\)
  & \((2.2686520413738664\times10^{-7},\,1.1943543223104596\times10^{-8})\) \\
100 & \(-6.553175069257122\times10^{-7}\)
  & \((1.235704860413911\times10^{-5},\,3.52904263500875\times10^{-6})\)
  & \((1.2128187577154448\times10^{-5},\,1.192559771676947\times10^{-6})\) \\
\bottomrule
\end{tabular}%
}
\end{table}

There were no componentwise sign reversals.  The order-\(129\) versus
order-\(257\) value changes were \(4.71\times10^{-14}\),
\(2.34\times10^{-13}\), and \(2.39\times10^{-12}\) at the three horizons; the
largest adjacent score change was \(1.07\times10^{-13}\).  Thus the observed
candidate differences are not explained by that particular adjacent
refinement.  They nevertheless remain descriptive: adjacent quadrature orders
can share truncation bias, no rigorous dense-reference error bound was proved,
and no target-specific equivalence margin was fixed before the run.

\paragraph{Trusted GPU/XLA execution and failure recording.}
The full \(T=1000,K=23\) manual-score route was executed in TensorFlow 2.19.1
on one logical RTX 4080 SUPER GPU.  The logical device was configured before
initialization with a hard \(8192\) MiB limit; memory growth was disabled
because TensorFlow does not permit both policies simultaneously.  XLA was
enabled and all claim-bearing outputs were placed on \texttt{GPU:0}.  The
concrete graph contained three functional \texttt{StatelessWhile} nodes, and a
source audit found no reachable Python horizon loop.

All center and four FD endpoint charts were valid.  The center value and score
were
\[
  \mathcal L(\theta_0)=-2286.226010065916,\qquad
  g^{\rm manual}=(5.664478697508284,-2.56604922190062).
\]
The GPU FD relative error was \(6.037244233113808\times10^{-8}\), and warm
replay, including the score, was bitwise identical.  Relative to the bound CPU
artifact, the maximum objective absolute difference was
\(9.094947017729282\times10^{-13}\), and the maximum manual-score absolute
difference was \(1.4408030324375432\times10^{-11}\).  These CPU/GPU
differences are descriptive; no equivalence threshold was added after seeing
them.

TensorFlow reported an allocator peak of \(1{,}995{,}776\) bytes.  This field
does not include the CUDA context, loaded libraries, or driver allocations and
must not be interpreted as total device use.  The run took \(69.5032\) seconds
in one attempt; this single timing is engineering telemetry, not a performance
distribution or method ranking.  TensorFlow's TF32 setting was recorded as
enabled, but the candidate tensors were float64, so the artifact makes no TF32
accuracy or throughput claim.

After module initialization, the harness writes a fresh structured failure
artifact for catchable Python, TensorFlow, configuration, identity, or
allocation exceptions raised by the main evaluation before propagating the
error.  Failures before that handler is installed, including import/device
initialization failures, and native process aborts cannot be serialized by this
mechanism.  The absence of a failure artifact is therefore not itself proof
that no such failure occurred.  Successful completion, JSON parsing, the
process exit status, and the saved output together establish the recorded
successful run.

\paragraph{Evidence boundary.}
The authoritative machine-readable evidence is in the dated Actual-SV
overcomplete analytical-chart repair directory under the repository's
benchmark evidence tree.
The selected preparation has SHA-256
\begin{quote}\small\ttfamily
e7958ba79da7f584b5e761faa22a9ed4f\allowbreak
dc53cd90027b0c13389388e130c6a8f.
\end{quote}
Exact commands, environment, device configuration, source state, timings, and
artifact paths are collected in the corresponding dated terminal result under
the repository's plan directory.  It is the Actual-SV overcomplete analytical
chart repair result dated 2026-07-17.

The resulting ledger action is deliberately narrow.  The former Actual-SV
fixed-square item is closed only as an experimental Contract E--TP local-chart
and own-scalar engineering defect: the overcomplete route covers the declared
finite points and differentiates its own finite scalar.  The evidence does not
close canonical Contract E--Chol wiring, cross-model propagation, scientific
score equivalence, trajectory-region validity, HMC readiness, default
readiness, or leaderboard admission.

\subsection{Algorithmic form and complexity}

The proposed local step is:
\begin{algorithm}[H]
\caption{Proposed Contract E--TP local score-aware reset}
\label{alg:bf-ledh-score-aware-teacher-projection}
\begin{algorithmic}[1]
  \State Compute the current LEDH corrected-weight increment \(\Delta_t\) and
    normalized source weights; add \(\Delta_t\) to the scalar.
  \State Form structurally valid teacher candidates
    \(z_{ij}=F_\theta(x_i,\varepsilon_j)\) and normalized weights \(\pi_{ij}\).
  \State Evaluate the fixed feature program \(\psi_\theta(z_{ij})\) and reduce
    \(b_\theta=\sum_{ij}\pi_{ij}\psi_\theta(z_{ij})\).
  \State Use a prepared fixed active set \(A\), or a reviewed full-support
    positive chart, and solve \(\Phi_{A,\theta}q_\theta=b_\theta\).
  \State Veto nonfinite, rank-deficient, nonpositive, or excessive-residual
    charts; do not clip weights or drop constraints.
  \State Carry \(\{s_\ell,q_\ell\}\) to the next LEDH step and initialize its
    prior log weights with \(\log q_\ell\), not \(-\log m\).
  \State In reverse mode, differentiate teacher construction, feature
    reduction, the same linear solve, carried log weights, and every later
    finite LEDH operation.
\end{algorithmic}
\end{algorithm}

With \(N\) parents, \(N\) innovation points, and \(K\) retained features, dense
teacher reduction costs \(O(N^2K)\).  A square chart adds \(O(K^3)\), while the
KKT form adds \(O(mK^2+K^3)\).  The projection is therefore semi-analytical;
the accepted \(N^2\) teacher work, not a large per-step optimizer, is expected
to dominate.  Streaming can reduce teacher storage from \(O(N^2d)\) to blocks,
provided the selected anchors and feature reductions are reproduced exactly in
the pullback.

\subsection{Concrete two-dimensional LGSSM witness}

The reference experiment uses
\begin{align}
  x_t&=A_\theta x_{t-1}+\varepsilon_t,
    &\varepsilon_t&\sim\mathcal N(0,Q_\theta),\\
  y_t&=h^\top x_t+\nu_t,
    &\nu_t&\sim\mathcal N(0,R_\theta),
\end{align}
with
\begin{equation}
\label{eq:bf-ledh-score-aware-lgssm-numbers}
  A_\theta=\begin{bmatrix}\rho&0.18\\-0.12&0.83\end{bmatrix},\quad
  Q_\theta=e^{2\ell_q}
    \begin{bmatrix}1&0.22\\0.22&0.65\end{bmatrix},\quad
  h=\begin{bmatrix}1\\0.35\end{bmatrix},\quad
  R_\theta=e^{2\ell_r}.
\end{equation}
This witness isolates the teacher and compression reset.  It uses the bootstrap
linear-Gaussian transition and observation potential directly; it does not
execute an LEDH migration or PF--PF proposal correction.  It can therefore
validate the finite teacher-projection propositions and the carried-weight
score mechanism, but not the surrounding LEDH flow implementation.
At
\[
  \theta_0=(\rho,\ell_q,\ell_r)=(0.72,-1.05,-0.70),
\]
the prior is
\[
  x_0\sim\mathcal N\!\left(
    \begin{bmatrix}0.25\\-0.30\end{bmatrix},
    \begin{bmatrix}0.70&0.16\\0.16&0.45\end{bmatrix}
  \right),\qquad (y_1,y_2)=(0.40,-0.15).
\]
The realized matrices are
\[
  A=\begin{bmatrix}0.72&0.18\\-0.12&0.83\end{bmatrix},\quad
  Q=\begin{bmatrix}
    0.1224564283&0.0269404142\\
    0.0269404142&0.0795966784
  \end{bmatrix},\quad
  R=0.2465969639.
\]

An order-nine tensor Gauss--Hermite rule in two dimensions gives 81 prior
parents and 81 innovation points, hence a 6561-point teacher at time one.  The
teacher weights are the Gauss--Hermite product weights multiplied by
\(g_\theta(y_1\mid z_{ij})\) and normalized.  The seven retained features are
\begin{equation}
\label{eq:bf-ledh-score-aware-lgssm-features}
  \psi_\theta(x)
  =\left(1,x_1,x_2,x_1^2,x_1x_2,x_2^2,r_{2,\theta}(x)\right)^\top,
\end{equation}
where the linear-Gaussian predictive feature is available analytically:
\begin{equation}
  r_{2,\theta}(x)
  =\mathcal N\!\left(
    y_2;h^\top A_\theta x,
    h^\top Q_\theta h+R_\theta
  \right).
\end{equation}

A deterministic linear-program preparation at \(\theta_0\) produced a positive
basic feasible set with zero-based teacher indices
\[
  A=(108,221,2317,2402,2474,3942,4001).
\]
Those indices are then frozen; only the seven-by-seven feature solve is
differentiated.  The resulting support and weights are:
\begin{center}
\begin{tabular}{rrrr}
\toprule
\(\ell\) & \(s_{\ell,1}\) & \(s_{\ell,2}\) & \(q_\ell\)\\
\midrule
1 & -3.4768598774 & -3.5565129989 & 0.0013273172\\
2 & -2.2613978461 & -1.2129585132 & 0.0002961221\\
3 & -0.5385388272 & -1.9705700614 & 0.0319724624\\
4 & -0.4079200865 & -0.1434149819 & 0.0377858212\\
5 & -0.6440561084 &  0.3400861366 & 0.0393703609\\
6 &  1.3859662461 & -1.8303729322 & 0.0541647250\\
7 &  0.4195516400 & -0.0203591973 & 0.8350831912\\
\bottomrule
\end{tabular}
\end{center}
The scaled feature matrix has condition number \(84.2606\), and the minimum
weight is \(2.96122\times10^{-4}\); hence this center lies inside, rather than on
the boundary of, the positive fixed chart.

Table~\ref{tab:bf-ledh-score-aware-lgssm-witness} separates teacher accuracy
from compression accuracy.  The Kalman filter is the exact LGSSM oracle.  The
teacher is a finite quadrature approximation to it.  The student is proved and
observed to reproduce the finite teacher for the retained predictive feature.
\begin{table}[H]
\centering
\caption{Two-observation 2D LGSSM value and score at
\(\theta_0=(0.72,-1.05,-0.70)\).}
\label{tab:bf-ledh-score-aware-lgssm-witness}
\resizebox{\textwidth}{!}{%
\begin{tabular}{lrrrr}
\toprule
Method & log likelihood & \(\partial_\rho\) & \(\partial_{\ell_q}\) &
  \(\partial_{\ell_r}\)\\
\midrule
Kalman oracle
 & -1.637545116868 & -0.843840883733 & -0.405394763108 & -0.667711074507\\
6561-point teacher
 & -1.637542132400 & -0.843638383670 & -0.405408322446 & -0.667677325884\\
7-point student
 & -1.637542132400 & -0.843638383670 & -0.405408322446 & -0.667677325884\\
Teacher minus Kalman
 & \(2.98447\times10^{-6}\) & \(2.02500\times10^{-4}\) &
   \(-1.35593\times10^{-5}\) & \(3.37486\times10^{-5}\)\\
Student minus teacher
 & \(0\) & \(0\) & \(5.55\times10^{-17}\) &
   \(-1.11\times10^{-16}\)\\
\bottomrule
\end{tabular}
}
\end{table}

The current teacher increment is \(-0.952850708925\), and the retained next
increment is \(-0.684691423474\) for both teacher and student.  In float64, the
maximum feature-value residual is \(1.30\times10^{-15}\), the maximum
feature-tangent residual over all seven features and three parameters is
\(4.05\times10^{-15}\), and the largest autodiff versus centered finite-
difference residual is \(1.83\times10^{-10}\) with step \(10^{-5}\).

The executable witness and its structured JSON result are stored under the
chapter's benchmark and benchmark-artifact directories.  Their exact paths,
command, and source state are recorded in the dated experiment result note.
It intentionally hides CUDA and runs as a small TensorFlow float64 CPU
reference.  It is implementation evidence for this numerical construction,
not GPU/XLA or production evidence.

\subsection{Mathematical and engineering boundary}

The propositions prove a precise finite statement: a positive coreset exists at
a fixed parameter value; a positive nonsingular frozen chart is locally
differentiable; all retained feature totals and tangents match; and retaining
the exact next finite predictive contribution makes the next value and score
match the teacher.  They do not prove any of the following:
\begin{itemize}
  \item that a low-order or finite teacher equals the exact nonlinear filter;
  \item that first two moments plus one predictive feature determine the entire
    tangent filtering measure;
  \item that one active set stays positive and nonsingular over an HMC
    trajectory or posterior region;
  \item that selecting a new active set online is differentiable;
  \item that a 6561-to-7 compression remains accurate beyond the one retained
    future observation;
  \item that full-state Contract E affine anchors preserve the deterministic
    support of a degenerate NAWM model; or
  \item that this proposed route is HMC-ready, canonical, default-ready, or
    computationally feasible at NAWM scale.
\end{itemize}
The next scientific tests must therefore use, in order, a multi-step LGSSM
rollout with Kalman value and score at every step, a singular structural-linear
fixture whose deterministic completion is exact, and only then a small pruned
DSGE model.  Passing the numerical witness above establishes feasibility of the
mechanism, not readiness of the algorithm.

\section{Model-agnostic comparison of finite-filter scores}
\label{sec:bf-ledh-score-comparison-opg}

LGSSM experiments are diagnostic because their Kalman oracle exposes failures
that would be hidden in a nonlinear model.  An LGSSM-specific expected Fisher
matrix is nevertheless the wrong general score-comparison interface: the
intended particle-filter applications usually do not have an analytic Fisher
matrix.  This section therefore defines two diagnostics that require only a
reference score and its time-local predictive increments.  They do not change
the finite LEDH scalar, its derivative, or Contract E.  They change only how a
candidate score error is reported.

Fix the parameter coordinates before comparison.  In an HMC experiment these
should normally be the unconstrained HMC coordinates.  Let
\begin{equation}
  \ell_T(\theta)=\sum_{t=1}^T \ell_t(\theta),\qquad
  s_t=\nabla_\theta\ell_t(\theta),\qquad
  g=\sum_{t=1}^T s_t,
  \label{eq:bf-ledh-score-increments}
\end{equation}
where
\(\ell_t(\theta)=\log p_\theta(y_t\mid y_{1:t-1})\) for the declared
reference filter.  If \(\widehat g\) is a candidate score in the same
coordinates, write \(\Delta=\widehat g-g\).

\subsection{Global score-vector norm diagnostics}

The direct global relative error is
\begin{equation}
  E_{\mathrm{total}}
  =\frac{\lVert\Delta\rVert_2}{\lVert g\rVert_2}.
  \label{eq:bf-ledh-score-total-relative}
\end{equation}
It is meaningful when \(\lVert g\rVert_2>0\), and must be reported as
undefined rather than zero when the denominator is zero.  Unlike a maximum of
componentwise relative errors, it does not diverge merely because one realized
score component is close to zero.  It can, however, hide an inaccurate small
component behind a large component, and the denominator can be small because
predictive score increments cancel.  Report the absolute error
\(\lVert\Delta\rVert_2\) with it.

A cancellation-resistant companion is
\begin{equation}
  E_{\mathrm{energy}}
  =\frac{\lVert\Delta\rVert_2}
  {\left(\sum_{t=1}^T\lVert s_t\rVert_2^2\right)^{1/2}}.
  \label{eq:bf-ledh-score-energy-relative}
\end{equation}
This denominator is the square root of the trace of the realized summed OPG
matrix, equivalently \(\sqrt{T\operatorname{tr}(\overline S_T)}\).
It is also undefined if every reference increment is zero.  Neither
\eqref{eq:bf-ledh-score-total-relative} nor
\eqref{eq:bf-ledh-score-energy-relative} is invariant to an arbitrary
rescaling of the parameter coordinates.  Consequently the coordinate map is
part of the artifact identity; a candidate and its reference must use the same
map.

\subsection{Average predictive-score OPG metric}

Define the realized average outer-product-of-gradients matrix
\begin{equation}
  \overline S_T=\frac{1}{T}\sum_{t=1}^T s_t s_t^\top.
  \label{eq:bf-ledh-average-opg}
\end{equation}
This is not \(gg^\top\).  The latter has rank at most one and repeats the
near-zero-denominator problem rather than estimating independent parameter
scales.  Nor is \(\overline S_T\) the outer product or covariance of scores
from repeated particle seeds on one observed data set.  Centered variation
across those seeds estimates Monte Carlo error; its uncentered second moment
mixes the observed score, approximation bias, and Monte Carlo noise.  Particle-
seed variation belongs in uncertainty intervals, not in the normalization
matrix.

For a diagonal-shrinkage coefficient \(0\leq\lambda\leq1\), a positive
diagonal scale matrix \(D\), a base ridge \(\epsilon_0>0\), and an optional
nonnegative numerical floor \(\epsilon_{\min}\), set
\begin{align}
  \rho_T&=\max\left\{\frac{\epsilon_0}{T},\epsilon_{\min}\right\},
  \label{eq:bf-ledh-average-opg-ridge}\\
  M_T&=(1-\lambda)\overline S_T
       +\lambda\operatorname{Diag}(\operatorname{diag}\overline S_T)
       +\rho_TD,
  \label{eq:bf-ledh-average-opg-metric}\\
  J_T&=T M_T.
  \label{eq:bf-ledh-total-opg-metric}
\end{align}
The matrix \(D\) supplies score-squared units and must be frozen independently
of the candidate error.  In dimensionless unconstrained coordinates, \(D=I\)
is a possible declared baseline, not a universal scientific default.  A
reference-derived diagonal scale may instead be frozen on separate simulated
trajectories or a tuning split.  Choosing \(D\), \(\lambda\), or either ridge
after inspecting which setting makes a candidate pass is forbidden.

The total-score quadratic diagnostic and its diagonal companion are
\begin{align}
  E_{\mathrm{OPG}}
  &=\left(\frac{1}{p}\Delta^\top J_T^{-1}\Delta\right)^{1/2},
  \label{eq:bf-ledh-score-opg-error}\\
  E_{\max,\mathrm{diag}}
  &=\max_{1\leq j\leq p}
    \frac{|\Delta_j|}{\sqrt{(J_T)_{jj}}}.
  \label{eq:bf-ledh-score-opg-diagonal-error}
\end{align}
The first detects errors in correlated directions and the second prevents a
small coordinate from being hidden by the global norm.  They are diagnostics,
not acceptance thresholds until a model-class-specific plan freezes a
criterion before candidate results are inspected.

\begin{proposition}[Finite-horizon well-posedness]
\label{prop:bf-ledh-average-opg-positive}
Suppose \(D\succ0\), \(\epsilon_0>0\), \(T<\infty\), and
\(0\leq\lambda\leq1\).  Then \(M_T\succ0\), even when \(T<p\), and
\eqref{eq:bf-ledh-score-opg-error} is finite for every finite \(\Delta\).
\end{proposition}

\begin{proof}
For every nonzero \(v\),
\begin{align*}
 v^\top M_Tv
 &=\frac{1-\lambda}{T}\sum_{t=1}^T(v^\top s_t)^2
   +\frac{\lambda}{T}\sum_{j=1}^p v_j^2
      \sum_{t=1}^T s_{t,j}^2
   +\rho_T v^\top Dv.
\end{align*}
The first two terms are nonnegative.  The last is strictly positive because
\(\rho_T\geq\epsilon_0/T>0\) and \(D\succ0\).  Hence \(M_T\), and therefore
\(J_T=T M_T\), is positive definite.
\end{proof}

The proposition makes a short-horizon diagnostic computable; it does not turn
the missing directions into observed information.  In particular,
\(\operatorname{rank}(\overline S_T)\leq\min(T,p)\).  For the current
\(T=2,p=5\) LGSSM diagnostic, at least three directions are supplied entirely
by shrinkage and ridge.  The artifact must report the unregularized rank or
eigenvalues, \(\rho_T\), and whether the numerical floor is active.

\begin{proposition}[Meaning of the \(1/T\) ridge]
\label{prop:bf-ledh-average-opg-vanishing-ridge}
Assume \(\epsilon_{\min}=0\), \(D_T\) is bounded, and
\(\overline S_T\to S\).  If \(\lambda_T\to\lambda_\star\), then
\[
  M_T\longrightarrow
  (1-\lambda_\star)S
  +\lambda_\star\operatorname{Diag}(\operatorname{diag}S),
\]
while
\[
  J_T
  =T\left[(1-\lambda_T)\overline S_T
  +\lambda_T\operatorname{Diag}(\operatorname{diag}\overline S_T)\right]
  +\epsilon_0D_T.
\]
Thus the ridge has constant size in the total-score metric and relative order
\(1/T\) in every persistently identified direction.
\end{proposition}

\begin{proof}
Substitute \(\rho_T=\epsilon_0/T\) into
\eqref{eq:bf-ledh-average-opg-metric} and multiply by \(T\).  Boundedness of
\(D_T\) makes \((\epsilon_0/T)D_T\to0\) in the average metric.  The stated
limits then follow by continuity of the diagonal map.
\end{proof}

Applying \(\epsilon_0/T\) directly to the summed OPG would make the ridge only
order \(1/T^2\) relative to its data term.  The average formulation is what
gives the intended order \(1/T\).  A fixed \(\lambda>0\) does not disappear
asymptotically: it converges to a diagonally shrunk metric, not to \(S\).
Recovering the unshrunk OPG requires \(\lambda_T\to0\).  Similarly, when
\(\epsilon_{\min}>0\) is active, the ridge no longer vanishes relative to the
average metric.  Persistent floor activity is evidence of weak excitation or
numerical scaling, and must be reported rather than hidden as successful
normalization.

Under the usual differentiability and interchange conditions at the true
parameter,
\(\mathbb E[s_t\mid y_{1:t-1}]=0\).  Predictive score increments are then a
martingale-difference sequence, so cross-time score covariances vanish and the
expected sum of \(s_ts_t^\top\) is the Fisher information of the observed-data
score.  On one realized trajectory, however, \(\overline S_T\) remains a
data-dependent OPG estimate.  For short or weakly excited records it should be
pooled over independently simulated or observed trajectories fixed before the
candidate comparison.  Exact LGSSM Fisher information is useful only as an
oracle check that this generic construction behaves sensibly; it is not the
normalization required by the nonlinear-model interface.

\section{Implementation contract}
\label{sec:bf-ledh-ot-implementation-contract}

A robust canonical implementation should expose two separable functions:
\[
  \operatorname{reset\_forward}(X,\log w,\Xi,\lambda;\eta_{\rm OT})
  =
  (X^\star,\mathrm{aux}),
  \qquad
  \operatorname{reset\_vjp}(\mathrm{aux},G_{X^\star})
  =
  (G_X,G_{\log w}),
\]
where \(\eta_{\rm OT}\) contains the declared OT settings: cost convention, epsilon,
annealing schedule, iteration budget, row and column chunk sizes, stabilization,
and stop-gradient policy.  The auxiliary record should contain only mathematical
or numerical data needed to reproduce the same finite map.  It should not
silently change the scalar by switching from streaming to dense semantics, from
finite Sinkhorn to another optimizer, or from one random stream to another.

The canonical TensorFlow implementation is owned by
\texttt{ledh\_contract\_e\_reset\_tf.py},
\texttt{ledh\_contract\_e\_streaming\_tf.py}, and the canonical finite filter
callable.  The streaming layer carries the augmented payload \([X,1]\), so its
forward and reverse programs retain both \(Q\) and \(M\) without materializing a
dense production matrix.  Its VJP exposes source-particle direct and transport
terms separately, probability-weight moment terms, normalized-log-weight
transport terms, and their correctly composed totals.  The checked LGSSM
callable packages a primal traversal and a separate five-direction manual-JVP
traversal; pointwise tests match the manual JVP to forward autodiff of the
private primal on frozen tiny fixtures.  This is not a claim that the manual
traversal is mechanically generated from one primal trace or cannot drift.

The following checks belong with any implementation:
\begin{enumerate}[label=(\arabic*)]
  \item row and column residuals for the finite transport object;
  \item equality of dense and streaming VJP on small problems where both are
  feasible;
  \item finite-difference or raw-autodiff agreement on small fixed-seed cases;
  \item explicit tests for stopped versus differentiated scale, keys, and
  potentials;
  \item static-shape behavior under the intended compiled path;
  \item finite value and finite score on the target model class before any HMC
  use.
\end{enumerate}

These checks certify local derivative accounting for the declared finite
computation.  They do not certify the statistical accuracy of the relaxed
filter, posterior convergence, or default production readiness.

\section{Relation to forward-mode JVP}
\label{sec:bf-ledh-ot-jvp}

Forward-mode JVP is also mathematically legitimate.  For a perturbation
\(\dot\theta\), it propagates tangents through the same filtering loop and
returns
\[
  D_\theta\mathcal L_K(\theta)\,\dot\theta .
\]
This can be memory efficient, and it is useful for diagnostics along a small
number of directions.  For HMC, however, the sampler usually needs the whole
score vector.  If the parameter dimension is \(d_\theta\), forward mode needs
roughly \(d_\theta\) directional passes unless it is batched over tangent
directions.  A VJP gives the full scalar-target score in one reverse pass.  That
is why the custom-gradient design in this chapter emphasizes the OT VJP.

\section{What this chapter does and does not claim}
\label{sec:bf-ledh-ot-boundary}

This chapter gives a mathematical route for a custom gradient of the executed
finite LEDH-PFPF-OT scalar.  For the canonical Contract E route it claims:
\begin{itemize}
  \item the raw barycentric OT product has the local VJP in
  \eqref{eq:bf-ledh-ot-vjp-barycentric}, but that object is historical and is
  not the canonical reset;
  \item the canonical finite transport uses the row quotient in
  \eqref{eq:bf-ledh-contract-e-row-quotient} and both quotient cotangents;
  \item the canonical total pullback includes both direct weighted-moment and
  streaming-transport terms as in
  \eqref{eq:bf-ledh-contract-e-total-source-pullback};
  \item the normalized finite-transport matrix has the VJP in
  \eqref{eq:bf-ledh-ot-vjp-logw}--\eqref{eq:bf-ledh-ot-vjp-cost};
  \item finite Sinkhorn or annealed Sinkhorn can be differentiated by a reverse
  scan through the executed softmin program;
  \item composing those local VJPs through the filtering loop gives the
  derivative of the same executed scalar under fixed branch and randomness
  assumptions.
\end{itemize}

It does \emph{not} claim that the relaxed OT filter equals a categorical
particle filter, that finite Sinkhorn equals unregularized OT, that a custom VJP
improves Monte Carlo accuracy, or that the resulting score is automatically
HMC-ready for a structural model.  It also does not claim that the checked tiny
Contract E derivative certificates establish Kalman equivalence, full-time
numerical adequacy, nonlinear state-support validity, v2 admission, leaderboard
completeness, or default readiness.  Those are separate validation questions.
The Actual-SV overcomplete result in
Section~\ref{sec:bf-ledh-actual-sv-overcomplete-certification} is likewise an
experimental finite-program certificate only.  It repairs the declared local
fixed-square chart and own-scalar derivative failure; it is not evidence that
the separate canonical, HMC-region, scientific-equivalence, or leaderboard
gates have passed.

\FloatBarrier
